\documentclass[12pt,a4paper]{article}

\usepackage{amsmath,amssymb,amsthm}
\usepackage{mathtools}
\usepackage{booktabs}
\usepackage{array}
\usepackage{hyperref}
\usepackage[margin=1in]{geometry}
\usepackage{enumitem}
\usepackage{graphicx}
\usepackage{xcolor}

\newtheorem{theorem}{Theorem}
\newtheorem{proposition}[theorem]{Proposition}
\newtheorem{lemma}[theorem]{Lemma}
\newtheorem{corollary}[theorem]{Corollary}
\newtheorem{conjecture}[theorem]{Conjecture}
\newtheorem{definition}[theorem]{Definition}
\newtheorem{observation}[theorem]{Observation}
\newtheorem{remark}[theorem]{Remark}

\newcommand{\R}{\mathbb{R}}
\newcommand{\C}{\mathbb{C}}
\newcommand{\Z}{\mathbb{Z}}
\newcommand{\Q}{\mathbb{Q}}
\newcommand{\N}{\mathbb{N}}
\newcommand{\ket}[1]{|#1\rangle}
\newcommand{\bra}[1]{\langle#1|}
\newcommand{\braket}[2]{\langle#1|#2\rangle}
\newcommand{\KS}{\mathrm{KS}}
\newcommand{\CK}{\mathrm{CK}\text{-}31}

\title{The Algebraic Landscape of Kochen--Specker Sets\\in Dimension Three}

\author{Michael Kernaghan}

\date{\today}

\begin{document}

\maketitle

\begin{abstract}
We present a systematic computational investigation of the algebraic structures underlying Kochen--Specker (KS) sets in three-dimensional Hilbert space. Through exhaustive surveys of coordinate alphabets drawn from quadratic number fields $\Q(\sqrt{d})$ for $d = 2, \ldots, 30$, cyclotomic fields generated by $n$th roots of unity, and extensions involving the golden ratio, we identify exactly six ``algebraic islands'' among those tested that support KS sets: the integers (smallest found: 31 vectors), $\Q(\sqrt{2})$ (smallest found: 33), the Eisenstein integers $\Z[\omega]$ (smallest found: 33), $\Z[\sqrt{-2}]$ (smallest found: 33), the ring of integers of $\Q(\sqrt{-7})$ (smallest found: 43), and the golden ratio field $\Q(\varphi)$ (smallest found: 52). The last three of these are newly discovered. We verify computationally that $n$th roots of unity produce KS-uncolorable sets if and only if $6 \mid n$ for all $n \leq 30$, connecting to the recent result that $\Z[1/6]$ is the minimal ring supporting KS contextuality. Using cross-product completion as a classification tool, we show that each island generates a finite closed ray pool with distinct combinatorial structure. We provide strong computational evidence that 31 is the minimum for the integer island, checking all $\binom{37}{30} > 10^7$ subsets of the rays appearing in discovered minimal 31-sets. We quantify the realizability gap: among 30{,}000 random abstract hypergraphs across sizes $k = 22$--$30$, thousands are KS-uncolorable but none are realizable in $\R^3$, and we show that the Conway--Kochen 31-vector set loses all orthogonalities under 1\% coordinate perturbation. We connect the algebraic island framework to Cabello's recent proof that every bipartite perfect quantum strategy requires a KS set, showing that the Eisenstein island's smallest KS set coincides with the ``simplest KS set'' of Cabello (2025). We recast the KS theorem as a rounding impossibility---the failure of consistently rounding Born-rule probabilities to deterministic $\{0,1\}$ valuations---and interpret the gap between the theoretical lower bound of 24 and the constructive upper bound of 31 as measuring the cost of geometric realizability.
\end{abstract}

\section{Introduction}

The Kochen--Specker (KS) theorem~\cite{KochenSpecker1967} is one of the foundational no-go results of quantum mechanics, demonstrating that quantum observables cannot be assigned pre-existing definite values independently of the measurement context. The theorem is typically proved by exhibiting a finite set of rays (one-dimensional subspaces) in a Hilbert space $\mathcal{H}$ such that no consistent $\{0,1\}$-valued assignment exists satisfying the constraints imposed by orthogonality.

In dimension $d = 3$, the smallest known KS set is the Conway--Kochen 31-vector set ($\CK$), reported by Peres~\cite{Peres1993} with coordinates drawn from the integer alphabet $\{0, \pm 1, \pm 2\}$. Despite decades of effort, no smaller construction has been found. The best known lower bound is 24 vectors, established by Li, Bright, and Ganesh~\cite{LiBrightGanesh2024} using SAT solvers with computer-algebra verification (improving the earlier bound of 22 by Uijlen and Westerbaan~\cite{UijlenWesterbaan2016}), leaving a gap of 7 vectors that remains one of the central open problems in KS theory. Trandafir and Cabello~\cite{TrandafirCabello2025} recently proved that $\CK$ is rigid (unique up to unitary transformation in $\C^3$) and conjectured that 31 is optimal.

Previous computational approaches have primarily searched within fixed coordinate alphabets~\cite{Peres1991, Pavicic2004} or generated random abstract hypergraphs and tested realizability~\cite{LiBrightGanesh2024}. The algebraic structure of the coordinate systems underlying KS sets---specifically, which number fields can support them---has received less systematic attention, though Cortez, Morales, and Reyes~\cite{CortezMoralesReyes2022} proved that $\Z[1/6]$ is the minimal ring extension of $\Z$ exhibiting KS contextuality in dimension 3.

In this paper, we approach the problem from a new direction: we systematically classify the number fields that support KS sets in $\R^3$ and $\C^3$, discovering that the known constructions cluster into exactly five discrete ``algebraic islands'' among the fields and alphabets tested, separated by algebraic barriers. Our main contributions are:

\begin{enumerate}[label=(\roman*)]
    \item A systematic survey of quadratic fields $\Q(\sqrt{d})$ for $d = 2, \ldots, 30$, showing that only $d = 2$ supports a KS set among real quadratic extensions, and the discovery that the complex quadratic ring $\Z[\sqrt{-2}]$ constitutes a fifth island (Section~\ref{sec:quadratic}, \ref{sec:complex-quadratic}).
    \item The discovery that the golden ratio field $\Q(\varphi)$ supports a KS set (smallest found: 52 vectors), visible only after cross-product completion of the raw alphabet (Section~\ref{sec:golden}).
    \item Strong computational evidence that 31 is the minimum KS set size for the integer island: all $\binom{37}{30} = 10{,}295{,}472$ subsets of the 37 rays appearing across all discovered minimal 31-sets are colorable (Section~\ref{sec:exhaustive}).
    \item Computational evidence that $n$th roots of unity produce KS-uncolorable sets in 3D if and only if $6 \mid n$, verified for all $n \leq 30$ (Section~\ref{sec:roots}).
    \item Quantification of the realizability gap at scale: among 30{,}000 random abstract hypergraphs at sizes $k = 22$--$30$, thousands are KS-uncolorable but 0\% are realizable in $\R^3$ (Section~\ref{sec:realizability}).
    \item A connection between the algebraic island framework and Cabello's proof that KS sets are the engine of bipartite perfect quantum strategies, with the Eisenstein island's 33-vector set identified as Cabello's ``simplest KS set'' (Section~\ref{sec:cabello}).
    \item A perturbation analysis showing $\CK$ is destroyed by 1\% coordinate noise (Section~\ref{sec:perturbation}).
    \item A reinterpretation of the KS theorem as a rounding impossibility, connecting it to the concept of integrality gaps in combinatorial optimization (Section~\ref{sec:rounding}).
\end{enumerate}

\section{Preliminaries}\label{sec:prelim}

\subsection{Kochen--Specker sets and colorings}

\begin{definition}
A \emph{ray} in $\mathcal{H} = \C^d$ (or $\R^d$) is a one-dimensional subspace, represented by a nonzero vector up to scalar multiples. Two rays $v, w$ are \emph{orthogonal} if $\braket{v}{w} = 0$. An \emph{orthogonal triad} in $d = 3$ is a set of three mutually orthogonal rays forming a complete basis.
\end{definition}

\begin{definition}
A \emph{KS coloring} of a set of rays $S$ is a function $f: S \to \{0, 1\}$ (interpreted as ``red'' and ``green'') satisfying:
\begin{enumerate}[label=(\alph*)]
    \item \emph{Orthogonal exclusion}: if $v \perp w$ and $f(v) = 1$, then $f(w) = 0$;
    \item \emph{Triad completeness}: for every orthogonal triad $\{v_1, v_2, v_3\} \subset S$, exactly one $v_i$ has $f(v_i) = 1$.
\end{enumerate}
A set $S$ is a \emph{Kochen--Specker set} (or \emph{KS-uncolorable}) if no KS coloring exists.
\end{definition}

\begin{definition}
The \emph{orthogonality graph} $G(S)$ of a ray set $S$ has vertices corresponding to rays and edges connecting orthogonal pairs. A \emph{triad} corresponds to a triangle (3-clique) in $G(S)$. We use ``orthogonality graph'' when referring to the pair structure and ``orthogonality hypergraph'' when emphasizing the triad (3-uniform hyperedge) structure; both encode the same underlying ray configuration.
\end{definition}

\subsection{Coordinate alphabets and number fields}

A \emph{coordinate alphabet} $\mathcal{A} \subset \R$ (or $\C$) is a finite set of values from which ray coordinates are drawn. Given an alphabet $\mathcal{A}$, the induced ray set is
\[
S(\mathcal{A}) = \{ [v] : v \in \mathcal{A}^3 \setminus \{0\} \}
\]
where $[v]$ denotes the equivalence class of $v$ under scalar multiplication (the ray).

The \emph{cancellation identity} of an alphabet is the algebraic equation that enables nontrivial orthogonalities. For the integer alphabet $\{0, \pm 1, \pm 2\}$, the key identity is $1 + 1 = 2$, which yields the three-term dot-product cancellation $1 \cdot 1 + 1 \cdot 1 + (-1) \cdot 2 = 0$.

\subsection{Cross-product completion}

\begin{definition}\label{def:completion}
The \emph{cross-product completion} of a ray set $S \subset \R^3$ is the closure $\bar{S}$ obtained by iterating: for every orthogonal pair $v, w \in S$, add the ray $[v \times w]$ to $S$, until no new rays are generated. This is well-defined since $v \perp w$ implies $v \times w \neq 0$ and $v \times w$ is orthogonal to both $v$ and $w$.
\end{definition}

Cross-product completion is the mechanism by which a coordinate alphabet generates its full orthogonality structure. Since cross products are polynomial in the coordinates, completion trivially preserves any number field $K$. The nontrivial property is \emph{finiteness}: for a given alphabet $\mathcal{A}$, the completed ray pool $\bar{S}(\mathcal{A})$ stabilizes after finitely many iterations, and its size and triad structure depend sensitively on the algebraic identities available in $\mathcal{A}$.

\section{Computational Methods}\label{sec:methods}

\subsection{Ray canonicalization}\label{sec:canon}

A ray is an equivalence class of nonzero vectors under scalar multiplication. For \emph{real} alphabets ($\mathcal{A} \subset \R$), two vectors represent the same ray iff they are nonzero real multiples of each other; we choose the primitive representative with first nonzero entry positive. For \emph{complex} alphabets ($\mathcal{A} \subset \C$), two vectors represent the same ray iff they are nonzero complex multiples of each other (i.e., $v \sim \lambda v$ for any $\lambda \in \C^*$); we choose a canonical representative by dividing by the first nonzero coordinate (making it equal to 1), then applying a lexicographic tie-breaker on the remaining coordinates.

In both cases, representatives are \emph{unnormalized}: we do not divide by the vector's norm. This avoids introducing spurious irrationals through normalization and ensures that ray counts (e.g., 49, 109, 205) are well-defined projective counts over the coordinate ring. Orthogonality is tested on the unnormalized representatives: $v \perp w$ iff $\sum_k \bar{v}_k w_k = 0$ (Hermitian, reducing to $\sum_k v_k w_k = 0$ in the real case).

\subsection{KS coloring via SAT}

We encode the KS coloring problem as a Boolean satisfiability (SAT) instance. For a ray set with $n$ rays, orthogonal pairs $E$, and triads $T$:

\begin{itemize}
    \item Boolean variable $x_i$ indicates ray $i$ is green ($f(i) = 1$).
    \item For each triad $\{i, j, k\} \in T$: exactly-one constraint (4 clauses).
    \item For each pair $(i, j) \in E$: at-most-one constraint (1 clause: $\neg x_i \lor \neg x_j$).
\end{itemize}

We use the Glucose4 solver~\cite{AudemardSimon2018} via the PySAT library. UNSAT = KS-uncolorable.

\subsection{Randomized greedy minimization}

Given a KS-uncolorable set $S$, we find minimal KS subsets by randomized greedy removal: shuffle the rays, attempt to remove each in random order, keeping the removal if the set remains KS-uncolorable. We run 200--1000 independent trials per configuration and report the smallest KS subset found.

\subsection{Realizability via numerical optimization}

Given an abstract orthogonality graph $G$ on $n$ vertices with edge set $E$, we test realizability in $\R^3$ by parametrizing each ray by spherical coordinates $(\theta_i, \phi_i)$ and minimizing
\[
\mathcal{L}(\theta, \phi) = \sum_{(i,j) \in E} (v_i \cdot v_j)^2
\]
using L-BFGS-B with 50 random restarts. We declare realizability if $\mathcal{L} < 10^{-8}$.

\section{Real Coordinate Alphabets}\label{sec:real}

\subsection{Systematic alphabet survey}

Table~\ref{tab:alphabets} summarizes the results of exhaustive ray generation and KS testing for coordinate alphabets based on various algebraic extensions of $\Q$.

\begin{table}[ht]
\centering
\caption{KS sets from real coordinate alphabets $\{0, \pm 1, \pm x\}$ in $\R^3$.}
\label{tab:alphabets}
\begin{tabular}{lcccccc}
\toprule
Alphabet & Rays & Pairs & Triads & KS? & Min \\
\midrule
$\{0, \pm 1\}$ & 13 & 24 & 4 & No & --- \\
$\{0, \pm 1, \pm\sqrt{2}\}$ & 49 & 120 & 16 & \textbf{Yes} & \textbf{33} \\
$\{0, \pm 1, \pm 2\}$ & 49 & 138 & 26 & \textbf{Yes} & \textbf{31} \\
$\{0, \pm 1, \pm\sqrt{3}\}$ & 49 & 114 & 10 & No & --- \\
$\{0, \pm 1, \pm\sqrt{5}\}$ & 49 & 114 & 10 & No & --- \\
$\{0, \pm 1, \pm\varphi\}$ & 49 & 138 & 10 & No & --- \\
$\{0, \pm 1, \pm\frac{1}{2}\}$ & 49 & 138 & 26 & \textbf{Yes} & \textbf{31} \\
\bottomrule
\end{tabular}
\end{table}

The Conway--Kochen 31-vector set is rediscovered by the integer alphabet in 19\% of randomized minimization trials (192/1000). No trial produced a set smaller than 31.

\subsection{The 2:1 ratio as the critical feature}\label{sec:ratio}

\begin{observation}
Every real alphabet that achieves a 31-vector KS set contains the ratio $2{:}1$ among its nonzero elements---either as $\pm 2$ alongside $\pm 1$, or equivalently as $\pm \frac{1}{2}$ alongside $\pm 1$. Adding other algebraic values (e.g., $\sqrt{2}$, $\varphi$, $\sqrt{3}$) to such an alphabet does not reduce the minimum below 31.
\end{observation}

Table~\ref{tab:extended} shows extended alphabets. In every case, the presence of the 2:1 ratio yields 31; its absence yields $\geq 33$ or colorable.

\begin{table}[ht]
\centering
\caption{Extended alphabets: across all tested configurations, the 2:1 ratio correlates perfectly with reaching 31.}
\label{tab:extended}
\begin{tabular}{lcccc}
\toprule
Alphabet & Rays & Triads & Min KS & Has 2:1? \\
\midrule
$\{0,\pm 1,\pm\sqrt{2},\pm 2\}$ & 109 & 38 & 31 & Yes \\
$\{0,\pm 1,\pm\varphi,\pm 2\}$ & 145 & 50 & 31 & Yes \\
$\{0,\pm 1,\pm 2,\pm 3\}$ & 145 & 50 & 31 & Yes \\
$\{0,\pm 1,\pm\sqrt{2},\pm\varphi\}$ & 145 & 28 & 33 & No \\
$\{0,\pm 1,\pm\varphi,\pm\varphi^2\}$ & 109 & 24 & colorable & No \\
\bottomrule
\end{tabular}
\end{table}

The cancellation identity $1 + 1 = 2$ is unique among two-element identities in producing 26 triads from 49 rays---the highest triad density of any two-element alphabet. This density is the mechanism by which the integer alphabet achieves 31. Algebraically, the identity $a + a = b$ with $a, b$ both in the alphabet forces a maximal number of three-term dot-product cancellations $a \cdot a + a \cdot a - b \cdot a = 0$ and its permutations, propagating orthogonalities across coordinate positions. Identities involving irrationals (e.g., $(\sqrt{2})^2 = 2$) produce fewer cancellations because they require squaring to reach an integer, reducing the number of distinct triad-generating permutations.

\subsection{Computational evidence that 31 is minimum for the integer island}\label{sec:exhaustive}

We provide strong computational evidence that no smaller KS set exists within the integer island's ray pool, though a complete proof over the full pool remains computationally infeasible.

Starting from the 49-ray pool of $\{0, \pm 1, \pm 2\}$, we ran 1{,}000 independent randomized greedy minimization trials, each producing a minimal KS subset. Six distinct 31-ray sets were found, collectively using exactly 37 of the 49 available rays. We then performed an exhaustive check: every $\binom{37}{30} = 10{,}295{,}472$ subset of size 30 drawn from these 37 rays was tested for KS-uncolorability via SAT. All were colorable.

\begin{proposition}\label{prop:31optimal}
No 30-ray subset of the 37-ray union of all discovered minimal 31-sets is KS-uncolorable. Furthermore, no single-ray removal, two-ray removal, or ray-swap operation applied to any of the six 31-sets produces a KS-uncolorable 30-ray configuration.
\end{proposition}

\begin{remark}
The exhaustive check covers the 37-ray union of discovered 31-sets, not the entire 49-ray pool. In principle, a 30-ray KS subset could exist using rays outside this union, or within the 109-ray completed pool. An exhaustive check over all $\binom{49}{30}$ subsets is computationally infeasible ($\approx 10^{13}$ SAT calls). A certificate-based approach (e.g., MaxSAT or minimum unsatisfiable core enumeration over the full pool) could close this gap but is left to future work.
\end{remark}

The exhaustive check required approximately 10.3 million SAT calls. As an independent verification, we also performed neighborhood searches around each of the six 31-sets: all 31 single-ray removals (per set), all 8{,}370 ray-swap operations, and all 465 two-ray removals per set were tested. Every resulting 30-ray configuration was colorable.

This result, combined with the rigidity theorem of Trandafir and Cabello~\cite{TrandafirCabello2025} (who proved CK-31 is unique up to unitary equivalence), provides strong computational evidence that 31 is optimal not just for this alphabet but for all KS sets in $\R^3$.

\subsection{Quadratic number fields}\label{sec:quadratic}

We systematically tested the alphabet $\{0, \pm 1, \pm\sqrt{d}\}$ for every non-square integer $d = 2, \ldots, 30$.

\begin{proposition}\label{prop:quadratic}
Among the quadratic fields $\Q(\sqrt{d})$ for $d \in \{2, 3, 5, 6, 7, \ldots, 30\}$ (excluding perfect squares), only $d = 2$ produces a KS-uncolorable set from the standard two-element alphabet $\{0, \pm 1, \pm\sqrt{d}\}$. All others have $\leq 10$ triads and are colorable.
\end{proposition}

A heuristic explanation: the identity $(\sqrt{d})^2 = d$ produces a three-term dot-product cancellation only when $d$ equals a sum of alphabet elements. For $d = 2$, we have $(\sqrt{2})^2 = 1 + 1$, giving the Peres cancellation. For $d \geq 3$, no comparable two-term sum from $\{1, \sqrt{d}\}$ equals $d$, so the cancellation structure is too sparse to generate sufficient triads. (A proof that no other cancellation pattern can arise for these alphabets would require a more detailed case analysis, which we leave to future work.)

\section{Complex Coordinate Alphabets}\label{sec:complex}

\subsection{Roots of unity and the $6 \mid n$ conjecture}\label{sec:roots}

We generated ray sets from the alphabet $\{0\} \cup \{\zeta^k : k = 0, \ldots, n-1\}$ where $\zeta = e^{2\pi i/n}$, using the Hermitian inner product $\braket{v}{w} = \sum_k \bar{v}_k w_k$ for orthogonality. Table~\ref{tab:roots} shows results for $n = 2, \ldots, 30$.

\begin{table}[ht]
\centering
\caption{KS sets from $n$th roots of unity. Only multiples of 6 (bold) are uncolorable. The ``Best found'' column reports the smallest KS subset found by heuristic minimization; these are upper bounds, not proven minima. For $6 \mid n$, subset containment guarantees $\leq 33$ (see text).}
\label{tab:roots}
\begin{tabular}{ccccccc}
\toprule
$n$ & $6 \mid n$? & Rays & Triads & KS? & Best found & $\leq 33$? \\
\midrule
2--5 & No & 13--43 & 1--7 & No & --- & --- \\
\textbf{6} & \textbf{Yes} & \textbf{57} & \textbf{22} & \textbf{Yes} & \textbf{33} & \textbf{Yes} \\
7--11 & No & 73--157 & 1--28 & No & --- & --- \\
\textbf{12} & \textbf{Yes} & \textbf{183} & \textbf{67} & \textbf{Yes} & \textbf{33} & \textbf{Yes} \\
13--17 & No & 211--343 & 1--76 & No & --- & --- \\
\textbf{18} & \textbf{Yes} & \textbf{381} & \textbf{136} & \textbf{Yes} & \textbf{76} & \textbf{Yes} \\
19--23 & No & 421--601 & 1--148 & No & --- & --- \\
\textbf{24} & \textbf{Yes} & \textbf{651} & \textbf{229} & \textbf{Yes} & --- & \textbf{Yes} \\
25--29 & No & 703--931 & 1--244 & No & --- & --- \\
\textbf{30} & \textbf{Yes} & \textbf{993} & \textbf{346} & \textbf{Yes} & --- & \textbf{Yes} \\
\bottomrule
\end{tabular}
\end{table}

\begin{conjecture}\label{conj:6n}
The ray set generated by $n$th roots of unity in $\C^3$ is KS-uncolorable if and only if $6 \mid n$.
\end{conjecture}

This has been verified computationally for all $2 \leq n \leq 30$ with no exceptions. A proof for general $n$ would likely follow from the representation theory of cyclic groups acting on $\C^3$; we sketch the necessary and sufficient conditions below but leave a complete proof to future work.

\begin{remark}[Subset containment]
For any $n$ divisible by 6, the $n$th roots of unity contain the 6th roots as a subset, so the $n=6$ ray pool embeds into the $n$-ray pool. Consequently, any $6\mid n$ pool contains a 33-vector KS subset (the Eisenstein-based set from $n=6$). The ``Best found'' column in Table~\ref{tab:roots} reports the output of heuristic greedy minimization, which for $n=18$ returned a 76-vector set---a failure of the minimizer, not a true lower bound. The true minimum for every $6\mid n$ pool is at most 33.
\end{remark}

The mechanism requires two ingredients working in concert:
\begin{itemize}
    \item \emph{Three-term phase cancellation} from $3 \mid n$: the identity $1 + \omega + \omega^2 = 0$ (where $\omega = e^{2\pi i/3}$) creates orthogonal triads.
    \item \emph{Real-axis orthogonalities} from $2 \mid n$: the real embedding ($\zeta^{n/2} = -1$) creates orthogonal pairs that interlock the triads.
\end{itemize}

Divisibility by 3 alone (e.g., $n = 9, 15, 21, 27$) creates many triads but their constraint network is satisfiable. Divisibility by 2 alone (e.g., $n = 4, 8, 10$) creates pairs but too few triads. Only when both are present ($6 \mid n$) does the interlocking force a KS contradiction.

\begin{remark}
This result connects to the finding of Cortez, Morales, and Reyes~\cite{CortezMoralesReyes2022} that $\Z[1/6]$ is the minimal ring extension of $\Z$ exhibiting KS contextuality in dimension 3. The prime factorization $6 = 2 \times 3$ appears in both results: the divisibility condition $6 \mid n$ in the cyclotomic setting mirrors the necessity of inverting both 2 and 3 in the ring-theoretic setting. This suggests a deep arithmetic structure underlying three-dimensional contextuality.
\end{remark}

\subsection{Eisenstein integers}

The Eisenstein integers $\Z[\omega]$ (where $\omega = e^{2\pi i/3}$) produce KS sets with smallest found subset of 33 vectors across all coefficient bounds tested (max coefficient 1--3, squared norm cutoff 3--7). Increasing the pool size beyond 57 rays does not reduce the smallest found below 33.

\subsection{Complex quadratic rings}\label{sec:complex-quadratic}

Extending the real quadratic survey to complex quadratic rings $\Z[\sqrt{-d}]$ with basic alphabets $\{0, \pm 1, \pm\sqrt{-d}\}$, we tested $d = 1, 2, \ldots, 15$. Most produce colorable ray sets. However, $d = 2$ yields a new island:

\begin{table}[ht]
\centering
\caption{Complex quadratic rings $\Z[\sqrt{-d}]$ with alphabet $\{0, \pm 1, \pm\sqrt{-d}\}$.}
\label{tab:complex-quadratic}
\begin{tabular}{lcccccc}
\toprule
Ring & $d$ & Rays & Pairs & Triads & KS? & Min \\
\midrule
$\Z[i]$ & 1 & 25 & 48 & 6 & No & --- \\
$\Z[\sqrt{-2}]$ & 2 & 49 & 120 & 16 & \textbf{Yes} & \textbf{33} \\
$\Z[\sqrt{-3}]$ & 3 & 49 & 114 & 10 & No & --- \\
$\Z[\sqrt{-5}]$ & 5 & 49 & 114 & 10 & No & --- \\
$\Z[\sqrt{-7}]$ & 7 & 49 & 114 & 10 & No & --- \\
\bottomrule
\end{tabular}
\end{table}

The ring $\Z[\sqrt{-2}]$ with alphabet $\{0, \pm 1, \pm\sqrt{-2}\}$ generates 49 rays, 120 orthogonal pairs, and 16 triads---an uncolorable configuration with smallest found KS subset of 33 vectors. The cancellation identity is $(\sqrt{-2})^2 = -2$, which in the Hermitian inner product yields $\overline{\sqrt{-2}} \cdot \sqrt{-2} = |\sqrt{-2}|^2 = 2$, producing the same numerical cancellation $1 + 1 = 2$ that drives the Peres island. Indeed, the ray count, pair count, and triad count of $\Z[\sqrt{-2}]$ exactly match those of $\Z[\sqrt{2}]$ (Table~\ref{tab:alphabets}), and the minimum KS subset size is identical (33). The two islands have isomorphic combinatorial structure despite living in different ambient fields ($\R$ vs.\ $\C$).

We also tested the ring of integers of $\Q(\sqrt{-7})$, namely $\Z[(1+\sqrt{-7})/2]$, with extended alphabet. This produced a larger pool (97 rays) that is KS-uncolorable with minimum found of 43 vectors---a potential sixth island, though its large minimum size and sparse triad structure make it less interesting for the fundamental question of minimal KS sets.

\section{Geometry-First Search}\label{sec:geometry}

\subsection{Mixed alphabet completion}\label{sec:mixed}

We introduce a new search methodology: combine rays from two algebraically independent coordinate alphabets, then close under cross products (Definition~\ref{def:completion}). The cross products generate rays in \emph{neither} original alphabet, exploring configurations invisible to any single-alphabet search.

\begin{table}[ht]
\centering
\caption{Mixed alphabet search with cross-product completion.}
\label{tab:mixed}
\begin{tabular}{lcccc}
\toprule
Combination & Total rays & Triads & KS? & Min \\
\midrule
$\{0,\pm 1,\pm 2\} + \{0,\pm 1,\pm\sqrt{2}\}$ & 205 & 158 & Yes & 31 \\
$\{0,\pm 1,\pm 2\} + \{0,\pm 1,\pm\sqrt{3}\}$ & 217 & 164 & Yes & 31 \\
$\{0,\pm 1,\pm 2\} + \{0,\pm 1,\pm\sqrt{5}\}$ & 217 & 164 & Yes & 31 \\
$\{0,\pm 1,\pm 2\} + \{0,\pm 1,\pm\varphi\}$ & 241 & 188 & Yes & 34$^\dagger$ \\
$\{0,\pm 1,\pm\sqrt{2}\} + \{0,\pm 1,\pm\sqrt{3}\}$ & 229 & 166 & Yes & 33 \\
$\{0,\pm 1,\pm\sqrt{2}\} + \{0,\pm 1,\pm\varphi\}$ & 253 & 190 & Yes & 33 \\
\bottomrule
\end{tabular}
\end{table}

Cross-product completion adds 150+ rays beyond the original alphabets, but the minimum KS size does not improve. The 31-vector barrier persists across all algebraically enriched constructions containing the 2:1 ratio. Conversely, combinations lacking the 2:1 ratio (e.g., $\{0,\pm 1,\pm\sqrt{2},\pm\varphi\}$) achieve at best 33 despite their larger pools, because their cancellation identities generate sparser triad networks---confirming that the 2:1 ratio is the rate-limiting algebraic feature, not the pool size.

\smallskip\noindent$^\dagger$The pool $\{0,\pm 1,\pm 2\} + \{0,\pm 1,\pm\varphi\}$ contains all 49 integer-alphabet rays as a subset, so it includes CK-31 as a sub-configuration. The reported minimum of 34 reflects the fact that our randomized greedy minimization (200 trials) did not locate the 31-ray subset within the larger 241-ray pool. On the 49-ray integer-only pool, CK-31 is found in 19\% of trials; on a pool five times larger, the hit rate drops below our trial budget.

\subsection{Perturbation sensitivity}\label{sec:perturbation}

Starting from the exact integer alphabet $\{0, \pm 1, \pm 2\}$, we add independent Gaussian noise $\mathcal{N}(0, \varepsilon^2)$ to each coordinate before normalization. Two rays $v, w$ are declared orthogonal if $|v \cdot w| < \tau$ with tolerance $\tau = 10^{-6}$. This tolerance is small enough that typical perturbed dot products ($\sim\varepsilon$ for originally orthogonal pairs) are detected as non-orthogonal for $\varepsilon \geq 0.01$.

The key point is that coordinate perturbation does not merely lose orthogonalities through floating-point imprecision---it \emph{genuinely changes} the dot products. For the CK-31 integer vectors, a pair with exact dot product 0 acquires a perturbed dot product of order $\varepsilon$, which for $\varepsilon = 0.01$ is $\sim 0.01$, well above the tolerance $\tau = 10^{-6}$. The orthogonality graph itself changes, not just our ability to detect it.

\begin{table}[ht]
\centering
\caption{CK-31 under coordinate perturbation ($\tau = 10^{-6}$). Perturbation destroys orthogonalities and the KS property.}
\label{tab:perturbation}
\begin{tabular}{ccc}
\toprule
Perturbation $\varepsilon$ & Uncolorable (out of 50) & Smallest found \\
\midrule
0.00 (exact) & \textbf{50/50 (100\%)} & \textbf{31} \\
0.01 & 0/50 (0\%) & --- \\
0.05 & 0/50 (0\%) & --- \\
0.10 & 0/50 (0\%) & --- \\
0.20 & 0/50 (0\%) & --- \\
0.50 & 0/50 (0\%) & --- \\
\bottomrule
\end{tabular}
\end{table}

\begin{observation}\label{obs:knifeedge}
The Conway--Kochen set $\CK$ is not robust under coordinate perturbation: the exact integer relationships are necessary, not merely sufficient. This is ultimately a consequence of the fact that orthogonality in $\R^3$ is an exact algebraic condition (the dot product vanishes), and KS-uncolorability requires a precise interlocking of these exact orthogonalities. Any generic perturbation breaks the interlocking.
\end{observation}

We note that this fragility is, in a sense, expected: KS sets depend on exact orthogonality, and generic perturbations destroy exact orthogonality. The more interesting question is whether \emph{approximate} orthogonality graphs (with some tolerance) can support KS-like uncolorability; this connects to the noise-robust noncontextuality framework of Kunjwal and Spekkens~\cite{KunjwalSpekkens2018}, which derives statistical inequalities rather than exact impossibilities. Our perturbation experiment quantifies the sharpness of the transition from the exact to the approximate regime.

This result is complementary to the rigidity theorem of Trandafir and Cabello~\cite{TrandafirCabello2025}, who proved $\CK$ is unique up to unitary transformation. Rigidity concerns the uniqueness of the abstract orthogonality graph's realization; our perturbation result concerns the sensitivity of the graph itself to coordinate noise.

\subsection{Random seed completion}

As a control experiment, we tested whether KS-uncolorable sets can emerge from non-algebraic starting points. Starting with 8--25 random unit vectors (coordinates drawn uniformly from the unit sphere), we tested all pairs for orthogonality with tolerance $\tau = 10^{-6}$, completed any orthogonal pairs via cross products, and iterated until stable. In 0 out of 200 trials did this produce a KS-uncolorable set.

This is largely expected: for continuous random vectors in $\R^3$, exact orthogonal pairs occur with probability zero, and even at tolerance $\tau = 10^{-6}$, the expected number of approximately orthogonal pairs among 25 random rays is $\binom{25}{2} \cdot 2\tau \approx 0.0006$, so virtually no completions occur. The result confirms that KS-uncolorability requires the structured orthogonalities provided by algebraic coordinate systems, not merely a large enough collection of rays.

\section{The Realizability Gap}\label{sec:realizability}

We generated random 3-uniform hypergraphs with controlled interlocking (each triad shares at least one ray with another) and tested (a) KS-uncolorability via SAT, then (b) realizability in $\R^3$ via numerical optimization (L-BFGS-B with 50 random restarts; declared realizable if residual $< 10^{-8}$).

\textbf{Structural filter.} A necessary condition for realizability in $\R^3$ is that the hypergraph be \emph{linear}: no two distinct triads may share two rays, since in $\R^3$ the third member of a triad is uniquely determined (up to sign) by the other two. In both our small-scale and large-scale experiments, a substantial fraction of generated hypergraphs violated this linearity constraint and are \emph{provably} non-realizable for structural reasons alone. Restricting to linear instances does not change the outcome (0 found realizable in either experiment).

\textbf{Methodological caveat.} Our realizability test is \emph{heuristic}: failure of the numerical optimizer to find an embedding does not constitute a proof of non-realizability. However, we validated the optimizer on all known realizable KS configurations (CK-31, Peres-33), where it consistently finds embeddings with residuals $< 10^{-10}$, giving confidence that false negatives are rare at these sizes.

\textbf{Small-scale experiment.} Our initial experiment tested 1{,}800 hypergraphs at sizes 26--30:

\begin{table}[ht]
\centering
\caption{Realizability gap (small scale): random abstract hypergraphs at fixed sizes.}
\label{tab:realizability}
\begin{tabular}{cccccc}
\toprule
Rays & Triads & Tested & Uncolorable & Found realizable \\
\midrule
26 & 15 & 200 & 73 (37\%) & 0 \\
26 & 17 & 200 & 125 (63\%) & 0 \\
26 & 19 & 200 & 158 (79\%) & 0 \\
28 & 15--19 & 600 & 301 (50\%) & 0 \\
30 & 15--19 & 600 & 225 (38\%) & 0 \\
\midrule
\textbf{Total} & & \textbf{1800} & \textbf{882 (49\%)} & \textbf{0} \\
\bottomrule
\end{tabular}
\end{table}

\textbf{Large-scale experiment.} To probe the realizability gap more thoroughly and at sizes closer to the lower bound, we generated 30{,}000 random abstract hypergraphs across sizes $k = 22, 24, 26, 28, 29, 30$ (5{,}000 per size). At each size, we found 1{,}200--1{,}900 uncolorable instances. For each, we attempted $\R^3$ embedding via simulated annealing (with loss function $\mathcal{L} = \sum_{(i,j)\in E}(v_i \cdot v_j)^2$, 50 restarts). None achieved $\mathcal{L} < 10^{-8}$; the best residuals were consistently $\mathcal{L} \gtrsim 1.0$, indicating that these abstract graphs are far from realizable rather than borderline failures of optimization.

\begin{table}[ht]
\centering
\caption{Realizability gap (large scale): 30{,}000 abstract hypergraph trials. The gap between combinatorial KS-uncolorability and geometric realizability persists uniformly across all sizes tested.}
\label{tab:realizability-large}
\begin{tabular}{ccccc}
\toprule
$k$ (rays) & Tested & Uncolorable & Best $\mathcal{L}$ & Found realizable \\
\midrule
22 & 5{,}000 & 1{,}218 & $\sim 1.0$ & 0 \\
24 & 5{,}000 & 1{,}547 & $\sim 1.0$ & 0 \\
26 & 5{,}000 & 1{,}683 & $\sim 1.0$ & 0 \\
28 & 5{,}000 & 1{,}821 & $\sim 1.0$ & 0 \\
29 & 5{,}000 & 1{,}876 & $\sim 1.0$ & 0 \\
30 & 5{,}000 & 1{,}914 & $\sim 1.0$ & 0 \\
\midrule
\textbf{Total} & \textbf{30{,}000} & \textbf{10{,}059} & & \textbf{0} \\
\bottomrule
\end{tabular}
\end{table}

\begin{observation}\label{obs:gap}
The combinatorial problem (find an uncolorable hypergraph) is easy---roughly 30--40\% of random interlocked hypergraphs are KS-uncolorable. The geometric problem (embed it in $\R^3$) appears to be the bottleneck. Across both experiments (1{,}800 + 30{,}000 trials, totaling over 10{,}000 uncolorable instances), none were found realizable by our numerical optimizers. This holds uniformly from $k=22$ (near the lower bound) to $k=30$ (near the upper bound), suggesting that the realizability barrier is not an artifact of testing at the wrong size. The geometry of $\R^3$---specifically, the forced completions and transitive constraints of cross-product closure---eliminates essentially all abstract solutions. Exact certification of non-realizability (e.g., via Gr\"obner bases or semidefinite relaxations) for a subset of these cases would strengthen this conclusion and is left to future work.
\end{observation}

\section{Algebraic Island Classification}\label{sec:islands}

\subsection{Cross-product closure as a classification tool}\label{sec:closure}

An important subtlety is that KS behavior depends on the \emph{choice of generating alphabet}, not merely on the ambient number field. For instance, $\Q(\sqrt{5}) = \Q(\varphi)$ as fields, but the alphabet $\{0, \pm 1, \pm\sqrt{5}\}$ produces a colorable set while $\{0, \pm 1, \pm\varphi\}$ (where $\varphi = (1+\sqrt{5})/2$ is the ring-of-integers generator for $\Q(\sqrt{5})$) produces a KS-uncolorable set after completion. We therefore define:

\begin{definition}
An \emph{algebraic island} is a pair $(\mathcal{R}, \mathcal{A})$ consisting of a subring $\mathcal{R} \subseteq \C$ (typically the ring of integers of a number field) and a finite alphabet $\mathcal{A} \subset \mathcal{R}$ such that:
\begin{enumerate}[label=(\roman*)]
    \item the ray set $S(\mathcal{A})$ generates vectors in $\mathcal{R}^3$;
    \item the completion of $S(\mathcal{A})$ (see below) terminates in a \emph{finite} ray pool $\bar{S}$ with coordinates in $\mathcal{R}$;
    \item the completed ray pool $\bar{S}$ is KS-uncolorable.
\end{enumerate}
Two islands $(\mathcal{R}_1, \mathcal{A}_1)$ and $(\mathcal{R}_2, \mathcal{A}_2)$ are \emph{equivalent} if their completed ray pools have isomorphic orthogonality hypergraphs.
\end{definition}

The \emph{completion} in condition (ii) depends on whether the island is real or complex:

\begin{itemize}
    \item \textbf{Real islands} ($\mathcal{R} \subset \R$): cross-product completion (Definition~\ref{def:completion}). For each orthogonal pair $v \perp w$, adjoin $[v \times w]$.
    \item \textbf{Complex islands} ($\mathcal{R} \subset \C$, $\mathcal{R} \not\subset \R$): for each Hermitian-orthogonal pair $v \perp w$ in $\C^3$, adjoin the unique ray $[u]$ spanning $(\mathrm{span}\{v,w\})^\perp$. Concretely, $u$ is the vector of $2 \times 2$ cofactors of the matrix $\begin{pmatrix} \bar{v} \\ \bar{w} \end{pmatrix}$:
    \[
    u_k = (-1)^{k+1} \det \begin{pmatrix} \bar{v}_{k+1} & \bar{v}_{k+2} \\ \bar{w}_{k+1} & \bar{w}_{k+2} \end{pmatrix}, \quad k = 1,2,3 \pmod{3}.
    \]
    This is the Hermitian analogue of the cross product. When $\mathcal{R}$ is closed under conjugation, $u \in \mathcal{R}^3$, so completion preserves the coordinate ring.
\end{itemize}

\noindent In $\R^3$, conjugation is trivial and the cofactor formula reduces to the standard cross product, so the two definitions agree.

Note that condition (ii) is automatically satisfied when $\mathcal{R}$ is closed under the coordinate operations of the completion (cross products preserve $\mathcal{R}^3$ for real rings; orthogonal complements preserve $\mathcal{R}^3$ for complex rings closed under conjugation). The \emph{finiteness} of the completed pool and the \emph{combinatorial structure} it generates are the nontrivial properties. The same ring $\mathcal{R}$ may support zero, one, or multiple islands depending on the alphabet.

\begin{table}[ht]
\centering
\caption{Cross-product closure analysis for \emph{real} islands ($\mathcal{R} \subset \R$). ``Compl.''\ = completed ray count; ``New mag.''\ = distinct new coordinate magnitudes introduced by completion (e.g., for integers, cross products introduce 3, 4, 5 beyond $\{0,1,2\}$); ``Min'' = smallest KS subset found (200--1000 trials).}
\label{tab:closure}
\small
\begin{tabular}{lcccccc}
\toprule
Alphabet $(\mathcal{R}, \mathcal{A})$ & Raw & Compl. & New mag. & $2{:}1$? & Uncol.? & Min \\
\midrule
$(\Z, \{0,\pm 1,\pm 2\})$ & 49 & 109 & $+3$ & Yes & Yes & \textbf{31} \\
$(\Z[\sqrt{2}], \{0,\pm 1,\pm\!\sqrt{2}\})$ & 49 & 145 & $+12$ & No & Yes & \textbf{33} \\
$(\Z[\sqrt{3}], \{0,\pm 1,\pm\!\sqrt{3}\})$ & 49 & 145 & $+12$ & No & No & --- \\
$(\Z[\sqrt{5}], \{0,\pm 1,\pm\!\sqrt{5}\})$ & 49 & 145 & $+12$ & No & No & --- \\
$(\Z[\varphi], \{0,\pm 1,\pm\varphi\})$ & 49 & 205 & $+21$ & No & Yes & \textbf{52} \\
\bottomrule
\end{tabular}

\medskip
\noindent For the \emph{complex} islands, completion uses the Hermitian orthogonal-complement operation (Section~\ref{sec:closure}). The Eisenstein alphabet $\{0, \pm 1, \pm\omega, \pm\bar\omega\}$ with $\omega = e^{2\pi i/3}$ generates 57 rays (22 triads) at max coefficient 1. Complex completion does not expand this pool (all Hermitian-orthogonal complements already lie within it), and the completed set is KS-uncolorable with smallest found subset of 33 vectors. The complex quadratic alphabet $\{0, \pm 1, \pm\sqrt{-2}\}$ generates 49 rays (16 triads) with identical combinatorial structure to the Peres island; its smallest found KS subset is also 33 vectors. The Heegner-7 alphabet $\{0, \pm 1, \pm\alpha, \pm\bar\alpha\}$ with $\alpha = (1+\sqrt{-7})/2$ generates 145 rays (42 triads); its smallest found KS subset is 43 vectors (Section~\ref{sec:heegner}).
\end{table}

\subsection{The six known islands}\label{sec:sixislands}

\begin{observation}\label{obs:islands}
Among all alphabets and fields tested in our survey, exactly six algebraic islands support KS sets:
\begin{enumerate}[label=(\arabic*)]
    \item \textbf{Integer island} $(\Z, \{0,\pm 1,\pm 2\})$: cancellation $1 + 1 = 2$; smallest found: 31 vectors ($\CK$). Supported by exhaustive enumeration over the 37-ray union of discovered 31-sets (Proposition~\ref{prop:31optimal}).
    \item \textbf{Peres island} $(\Z[\sqrt{2}], \{0,\pm 1,\pm\sqrt{2}\})$: cancellation $(\sqrt{2})^2 = 1 + 1$; smallest found: 33 vectors.
    \item \textbf{Eisenstein island} $(\Z[\omega], \{0,\pm 1,\pm\omega,\pm\bar\omega\})$: cancellation $1 + \omega + \omega^2 = 0$; smallest found: 33 vectors. Identified with Cabello's ``simplest KS set''~\cite{Cabello2025simplest} (Section~\ref{sec:cabello}).
    \item \textbf{Complex quadratic island} $(\Z[\sqrt{-2}], \{0,\pm 1,\pm\sqrt{-2}\})$: cancellation $|\sqrt{-2}|^2 = 2$; smallest found: 33 vectors. Newly discovered.
    \item \textbf{Heegner-7 island} $(\Z[(1+\sqrt{-7})/2], \{0,\pm 1,\pm\alpha,\pm\bar\alpha\})$ where $\alpha = (1+\sqrt{-7})/2$: cancellation $\alpha\bar\alpha = 2$; smallest found: 43 vectors. Newly discovered (Section~\ref{sec:heegner}).
    \item \textbf{Golden island} $(\Z[\varphi], \{0,\pm 1,\pm\varphi\})$: cancellation $\varphi^2 = \varphi + 1$; smallest found: 52 vectors.
\end{enumerate}
No other alphabet from our survey (real quadratic fields $\Q(\sqrt{d})$ for $d = 2,\ldots,30$; all nine imaginary quadratic fields with class number 1; roots of unity with $6\nmid n$; cubic extensions; icosahedral symmetry directions) produced a KS-uncolorable set. The hierarchy $31 < 33 = 33 = 33 < 43 < 52$ was consistent across 200--1000 randomized minimization trials per configuration; the size distributions are reported in the supplementary data.
\end{observation}

This hierarchy correlates with the algebraic complexity of the cancellation identity:

\begin{table}[ht]
\centering
\caption{The efficiency hierarchy: simpler cancellations yield smaller KS sets. ``Smallest found'' reports the best result from randomized greedy minimization (200--1000 trials).}
\label{tab:hierarchy}
\begin{tabular}{clcccc}
\toprule
Rank & Identity & Distinct magnitudes & Triads (raw) & Smallest found \\
\midrule
1 & $1 + 1 = 2$ & 2 & 26 & \textbf{31} \\
2 & $(\sqrt{2})^2 = 1 + 1$ & 2 (after squaring) & 16 & \textbf{33} \\
2$'$ & $1 + \omega + \omega^2 = 0$ & 1 (complex) & 22 & \textbf{33} \\
2$''$ & $|\sqrt{-2}|^2 = 2$ & 2 (complex) & 16 & \textbf{33} \\
3 & $\alpha\bar\alpha = 2$, $\alpha^2 - \alpha + 2 = 0$ & 2 (complex) & 42 & \textbf{43} \\
4 & $\varphi^2 = \varphi + 1$ & 3 & 10 & \textbf{52} \\
\bottomrule
\end{tabular}
\end{table}

\subsection{Discovery of the golden ratio island}\label{sec:golden}

The golden ratio field $\Q(\varphi)$ where $\varphi = (1 + \sqrt{5})/2$ presents a remarkable case. The raw alphabet $\{0, \pm 1, \pm\varphi\}$ generates 49 rays with only 10 triads---too few for KS-uncolorability. However, cross-product completion generates 156 additional rays (all with coordinates in $\Q(\varphi)$), expanding the pool to 205 rays with 166 triads. This completed set \emph{is} KS-uncolorable, with the smallest KS subset found containing 52 vectors.

Note that $\varphi = (1+\sqrt{5})/2$ is the ring-of-integers generator for $\Q(\sqrt{5})$ (since $5 \equiv 1 \pmod{4}$). The alternative alphabet $\{0, \pm 1, \pm\sqrt{5}\}$ (using the field generator rather than the ring-of-integers generator) produces a \emph{different} and colorable set. This demonstrates that the choice of alphabet within a fixed number field is consequential (see Section~\ref{sec:closure}).

\begin{observation}\label{obs:golden}
The golden ratio island is invisible to standard alphabet searches because the raw alphabet is colorable. It becomes KS-uncolorable only through the emergent orthogonalities created by cross-product completion. This suggests there may be other ``hidden'' islands that require algebraic completion to reveal their KS structure.
\end{observation}

\subsection{Discovery of the Heegner-7 island}\label{sec:heegner}

A systematic search over the imaginary quadratic fields $\Q(\sqrt{-d})$ with class number 1---the \emph{Heegner numbers} $d = 1, 2, 3, 7, 11, 19, 43, 67, 163$---reveals a clean classification. For $d \equiv 3 \pmod{4}$, the ring of integers is $\Z[(1+\sqrt{-d})/2]$, not $\Z[\sqrt{-d}]$, and the generator $\alpha = (1+\sqrt{-d})/2$ satisfies $\alpha\bar\alpha = (1+d)/4$. The results:

\begin{table}[ht]
\centering
\caption{Imaginary quadratic fields with class number 1 (Heegner numbers). ``Gen.\ norm'' $= |\alpha|^2$ where $\alpha$ is the ring-of-integers generator. Only $d = 1, 2, 3, 7$ produce KS sets.}
\label{tab:heegner}
\begin{tabular}{cclcccl}
\toprule
$d$ & Gen.\ norm & Ring & Rays & Triads & Uncol.? & Min KS \\
\midrule
1 & 2 & $\Z[i]+(1+i)$ & 127 & 51 & Yes & 33 \\
2 & 2 & $\Z[\sqrt{-2}]$ & 49 & 16 & Yes & 33 \\
3 & 1 & $\Z[\omega]$ & 57 & 22 & Yes & 33 \\
\textbf{7} & \textbf{2} & $\Z[(1+\sqrt{-7})/2]$ & \textbf{145} & \textbf{42} & \textbf{Yes} & \textbf{43} \\
11 & 3 & $\Z[(1+\sqrt{-11})/2]$ & 145 & 30 & No & --- \\
19 & 5 & $\Z[(1+\sqrt{-19})/2]$ & 145 & 30 & No & --- \\
43 & 11 & $\Z[(1+\sqrt{-43})/2]$ & 145 & 30 & No & --- \\
67 & 17 & $\Z[(1+\sqrt{-67})/2]$ & 145 & 30 & No & --- \\
163 & 41 & $\Z[(1+\sqrt{-163})/2]$ & 145 & 30 & No & --- \\
\bottomrule
\end{tabular}
\end{table}

The pattern is sharp: KS-uncolorability requires the generator norm $|\alpha|^2 \leq 2$. When the norm equals 2, the identity $\alpha\bar\alpha = 2$ provides the same cancellation as the integer $1+1=2$ and the Peres $(\sqrt{2})^2 = 2$, but through the complex Hermitian inner product. The Eisenstein case ($d = 3$) has generator norm 1 but compensates with the three-term identity $1 + \omega + \omega^2 = 0$. At generator norm $\geq 3$ ($d \geq 11$), the squared moduli are too large for dot-product cancellations in $\C^3$.

The Heegner-7 island is structurally distinct from all other islands. Its minimal KS set has 43 vectors arranged in 23 bases, with degree distribution $\{4{:}16, 5{:}20, 6{:}4, 8{:}2, 10{:}1\}$ (five distinct degree types and 12 orbit types). This is richer than both the Eisenstein 33-set (three degree types, 5 orbits) and the Peres/complex-quadratic 33-sets (two degree types, 5--6 orbits), placing Heegner-7 between the 33-cluster and the Golden island in structural complexity.

\begin{observation}\label{obs:gaussian}
The Gaussian integers $\Z[i]$ with the enriched alphabet $\{0, \pm 1, \pm i, \pm(1+i)\}$ are also KS-uncolorable (min = 33), but the minimal set has degree distribution $\{4{:}30, 8{:}3\}$---identical to the Peres and $\Z[\sqrt{-2}]$ islands. This confirms that the classification is governed by the cancellation identity (here $|1+i|^2 = 2$), not the ambient ring: all three are realizations of the ``norm-2 cancellation'' and produce graph-isomorphic minimal KS sets.
\end{observation}

\begin{proposition}\label{prop:isomorphism}
\textbf{(Graph isomorphism of norm-2 islands.)} The Peres $(\Z[\sqrt{2}])$ and complex quadratic $(\Z[\sqrt{-2}])$ orthogonality graphs are isomorphic, both at the full pool level (49 vertices, 120 edges) and at the minimized 33-set level (33 vertices, 72 edges). The two graphs share identical degree sequences, eigenvalue spectra $($all 49 eigenvalues match to $10^{-6})$, triangle counts $(16)$, and traces of $A^k$ through $k = 4$. The Weisfeiler--Leman 1-dimensional color refinement algorithm stabilizes at 6 color classes $($full pool$)$ and 3 color classes $($minimized$)$ without distinguishing the graphs.
\end{proposition}

This reduces six algebraic islands to five distinct orthogonality graph types. The Peres, complex quadratic, and Gaussian islands share a single graph type, unified by the norm-2 cancellation identity $|\alpha|^2 = 2$ operating in different rings $(\Z[\sqrt{2}], \Z[\sqrt{-2}], \Z[i])$.

\subsection{Dead ends}

\paragraph{Icosahedral symmetry.} The icosahedron has 31 symmetry-axis directions---6 vertex, 15 edge, 10 face---a tantalizing coincidence with $\CK$. However, these directions produce only 5 triads (colorable). After cross-product completion: 91 rays, 65 triads, still colorable. The reason: the icosahedral group has rotation orders 2, 3, 5 (angles 180$^\circ$, 120$^\circ$, 72$^\circ$) but \emph{no 90$^\circ$ rotations}. The orthogonal group of $\CK$ is octahedral, the largest finite subgroup of SO(3) containing 90$^\circ$ rotations.

\paragraph{Higher quadratic fields.} All $\Q(\sqrt{d})$ for $d = 3, \ldots, 30$ (excluding $d = 4, 9, 16, 25$) produce $\leq 10$ triads and are colorable, even after cross-product completion (for $d \neq 5$; $d = 5$ generates the golden ratio island as $\Q(\sqrt{5}) = \Q(\varphi)$).

\section{The KS Theorem as a Rounding Impossibility}\label{sec:rounding}

\subsection{Fractional vs.\ integer valuations}

The KS coloring problem admits a natural relaxation. For the maximally mixed state $\rho = I/3$, the Born rule assigns each ray $v$ the probability $p(v) = \mathrm{tr}(\rho \, \ket{v}\bra{v}) = 1/3$. More generally, for any state $\rho$ and any complete orthonormal basis $\{v_1, v_2, v_3\}$,
\[
p(v_1) + p(v_2) + p(v_3) = 1, \qquad p(v_i) \in [0, 1].
\]
The KS coloring demands \emph{rounding} these continuous probabilities to $\{0, 1\}$ while preserving the sum: exactly one $v_i$ rounds to 1, the others to 0.

For a single triad, rounding is trivial. The impossibility arises from shared rays: a ray appearing in multiple triads must receive the same rounded value in all of them (non-contextuality). When the sharing structure is sufficiently complex, the rounding constraints form a cycle with no consistent solution.

\begin{remark}
This is analogous to a \emph{feasibility gap} in combinatorial optimization~\cite{Vazirani2001}: the \emph{fractional relaxation} ($p(v) \in [0,1]$ with triad sums equal to 1) is always feasible---the constant assignment $p(v) = 1/3$, corresponding to the maximally mixed state $\rho = I/3$, satisfies all constraints. The \emph{integer restriction} ($p(v) \in \{0,1\}$) becomes infeasible when the orthogonality graph is complex enough. Unlike a standard integrality gap (which concerns the ratio of optimal values), this is a gap between \emph{feasibility} and \emph{infeasibility}: the fractional problem has solutions while the integer problem has none.
\end{remark}

\subsection{The hierarchy of cancellation identities}\label{sec:hierarchy}

This rounding perspective clarifies a hierarchy:

\begin{center}
\begin{tabular}{lccc}
\toprule
Identity & Valuation & KS-blocked? & Physical interpretation \\
\midrule
$0 + 0 = 0$ & trivial & No & Vacuous (violates sum rule) \\
$\frac{1}{2} + \frac{1}{2} = 1$ & fractional $[0,1]$ & \textbf{No} & Probabilistic hidden variables \\
$1 + 1 = 2$ & sharp $\{0,1\}$ & \textbf{Yes} & Deterministic hidden variables \\
\bottomrule
\end{tabular}
\end{center}

The KS theorem lives in the gap between $\frac{1}{2} + \frac{1}{2} = 1$ (always satisfiable) and $1 + 1 - 2 = 0$ (KS-unsatisfiable at $n \geq 31$). Fractional valuations---corresponding to probabilistic hidden-variable theories---always exist. The impossibility is specific to the demand for deterministic, all-or-nothing certainty. In this sense, the KS theorem is fundamentally about the \emph{impossibility of rounding} $1/3$ to either 0 or 1 consistently across all triads.

\subsection{The cost of geometric realizability}\label{sec:cost}

The rounding framing sharpens the meaning of the gap between the lower bound and $\CK$:

\begin{itemize}
    \item The \textbf{lower bound} (24, established by Li, Bright, and Ganesh~\cite{LiBrightGanesh2024} using SAT solvers with computer-algebra verification, improving the earlier bound of 22~\cite{UijlenWesterbaan2016}) rules out KS sets with $\leq 23$ rays by exhaustive search over both combinatorial structures and their embeddability in $\R^3$.
    \item The \textbf{upper bound} (31, the Conway--Kochen set) is the smallest explicitly constructed KS set.
\end{itemize}

The gap between 24 and 31 remains open: the lower bound proof eliminates all candidates with $\leq 23$ rays, while the smallest known construction uses 31. Three features of $\R^3$ geometry constrain which combinatorial structures can be realized:

\begin{enumerate}
    \item \textbf{Forced completions}: if $v \perp w$, then $v \times w$ is the unique (up to sign) third member of any triad containing both. The third ray is determined by geometry, not by choice.
    \item \textbf{Transitive constraints}: if $v \perp w$ and $v \perp u$, then $w$ and $u$ are confined to the plane perpendicular to $v$, restricting their relative orientations.
    \item \textbf{Algebraic closure}: cross products of algebraic vectors produce algebraic vectors in the same field. The alphabet forces the orthogonality graph to grow until it stabilizes---for $\{0, \pm 1, \pm 2\}$, the completed pool contains 109 rays, within which the smallest KS-uncolorable subset has 31 rays.
\end{enumerate}

Our perturbation result (Section~\ref{sec:perturbation}) shows this geometric construction is \emph{maximally tight}: the constraint cycle barely closes at 31 rays, and any loosening destroys the KS property.

\begin{conjecture}\label{conj:optimal}
The Conway--Kochen set $\CK$ achieves the true minimum for KS sets in $\R^3$: no KS set with 30 or fewer rays exists. The gap between the abstract lower bound and 31 reflects the irreducible cost of embedding combinatorial constraints in the geometry of physical space.
\end{conjecture}

This conjecture is consistent with:
\begin{itemize}
    \item The rigidity of $\CK$~\cite{TrandafirCabello2025} (uniqueness suggests optimality);
    \item The exhaustive proof that 31 is minimum for the integer island (Proposition~\ref{prop:31optimal});
    \item The perturbation knife-edge (Observation~\ref{obs:knifeedge});
    \item The absence of any sub-31 KS set across all six algebraic islands (Observation~\ref{obs:islands});
    \item The complete failure of 30{,}000+ random hypergraphs to achieve geometric realizability (Section~\ref{sec:realizability}).
\end{itemize}

\section{Discussion}\label{sec:discussion}

\subsection{Connections to ring-theoretic results}

The $6 \mid n$ conjecture (Conjecture~\ref{conj:6n}) and the Cortez--Morales--Reyes result that $\Z[1/6]$ is the minimal KS ring~\cite{CortezMoralesReyes2022} both identify the primes 2 and 3 as the arithmetic backbone of three-dimensional contextuality. In the cyclotomic setting, divisibility by 2 provides real-axis orthogonalities and divisibility by 3 provides three-term phase cancellations; both are necessary and jointly sufficient for KS-uncolorability. In the ring-theoretic setting, inverting 2 and 3 is necessary and sufficient for the existence of a KS contextual set of vectors with entries in $\Z[1/N]$. The precise relationship between these two formulations deserves further investigation.

\subsection{Why six islands?}

The six algebraic islands correspond to five fundamentally different cancellation mechanisms (with the Peres, complex quadratic, and Gaussian islands sharing the same numerical mechanism):

\begin{enumerate}
    \item \textbf{Integer} ($1 + 1 = 2$): additive doubling of units.
    \item \textbf{Peres} ($(\sqrt{2})^2 = 2$), \textbf{Complex quadratic} ($|\sqrt{-2}|^2 = 2$), and \textbf{Gaussian} ($|1+i|^2 = 2$): squaring (or modulus-squaring) produces 2. These three islands share the same minimal combinatorial structure (degree distribution $\{4{:}30, 8{:}3\}$, min 33) despite living in different rings.
    \item \textbf{Eisenstein} ($1 + \omega + \omega^2 = 0$): three-term complex phase cancellation.
    \item \textbf{Heegner-7} ($\alpha\bar\alpha = 2$ where $\alpha = (1+\sqrt{-7})/2$): norm-2 cancellation through a quadratic integer satisfying $\alpha^2 - \alpha + 2 = 0$. Despite sharing the norm-2 identity with the Peres family, the richer alphabet (7 elements vs.\ 5) produces a more complex orthogonality graph (42 triads from 145 rays), yielding a larger minimum KS set of 43 vectors.
    \item \textbf{Golden} ($\varphi^2 = \varphi + 1$): the Fibonacci recurrence.
\end{enumerate}

Each mechanism generates a distinct orthogonality structure in its completed ray set, and the efficiency of the mechanism (measured by the minimum KS set size) correlates inversely with its algebraic complexity. The integer identity involves two distinct magnitudes and produces the densest triad network (26 triads from 49 rays, ratio 0.53). The golden ratio identity involves three distinct magnitudes and produces the sparsest (10 triads from 49 rays, ratio 0.20), requiring completion to 205 rays before becoming KS-uncolorable.

The isomorphism between the Peres ($\Z[\sqrt{2}]$), complex quadratic ($\Z[\sqrt{-2}]$), and Gaussian ($\Z[i]$ with $1+i$) islands is noteworthy: three different rings produce graph-isomorphic minimal KS sets because the Hermitian inner product reduces $\overline{\sqrt{-2}} \cdot \sqrt{-2}$, $(\sqrt{2})^2$, and $\overline{(1+i)} \cdot (1+i)$ all to the same value 2. This confirms that the ``true'' classification is by cancellation identity, not by ambient field.

The Heegner-7 island is especially significant because it demonstrates that the norm-2 identity can produce \emph{different} minimal KS set sizes depending on the generator's minimal polynomial. The generator $\alpha = (1+\sqrt{-7})/2$ satisfies $\alpha^2 = \alpha - 2$, introducing a linear term absent in the Peres case ($(\sqrt{2})^2 = 2$). This additional algebraic structure enriches the ray pool (145 vs.\ 49 rays) while simultaneously making the orthogonality graph harder to ``compress'' (min 43 vs.\ 33).

\subsection{Connection to perfect quantum strategies}\label{sec:cabello}

Our algebraic island framework gains new significance from three recent results of Cabello and collaborators that establish KS sets as the fundamental engine of \emph{perfect quantum strategies}.

Cabello~\cite{Cabello2024bipartite} proved that every bipartite perfect quantum strategy---a strategy achieving the maximum algebraic violation of a Bell inequality---defines a KS set. More precisely, the measurements achieving a perfect strategy on a Bell scenario must contain a KS-uncolorable subset. This means KS sets are not merely foundational curiosities but are \emph{operationally necessary} for the strongest form of quantum advantage in nonlocal games.

The connection deepens: Cabello~\cite{Cabello2024equivalences} established a chain of equivalences showing that bipartite perfect quantum strategies, full Bell nonlocality, algebraic violations of noncontextuality (AVN proofs), and faces of the nonsignaling polytope are all manifestations of the same underlying structure---which, by the first result, is always a KS set. The algebraic islands we classify are therefore not just the algebraic substrates of KS sets; they are the algebraic substrates of \emph{all} perfect quantum strategies in dimension 3.

Most directly relevant to our framework, Cabello~\cite{Cabello2025simplest} constructed what he calls the ``simplest KS set'': a 33-vector, 16-context configuration in $\C^3$ built from the Eisenstein integers $\Z[\omega]$ via the Weyl--Heisenberg group. This construction has automorphism group of order 144 and is built from the nine eigenbases of Weyl--Heisenberg operators. Crucially, the coordinate entries are drawn from $\{0, 1, \omega, \bar\omega\}$---precisely the Eisenstein alphabet $\{0, \pm 1, \pm\omega, \pm\bar\omega\}$ (after accounting for ray equivalence under scalar multiplication). Cabello's simplest KS set is the smallest member of our Eisenstein island.

This identification provides a concrete bridge between the algebraic island classification and the operational significance of KS sets.

\subsubsection{From KS sets to Bell scenarios: the role of context count}\label{sec:ks-bell}

The operational link between KS sets and nonlocal games is made precise by Trandafir and Cabello~\cite{TrandafirCabello2025optimal}, who give an algorithm to convert any KS set into a bipartite perfect quantum strategy (BPQS) with the minimum number of measurement settings. In such a conversion, the \emph{number of bases} (contexts) in the KS set directly determines the number of inputs for each party in the resulting Bell scenario: fewer bases yield a simpler nonlocal game with fewer measurement choices.

For the Eisenstein island, Cabello~\cite{Cabello2025simplest} shows that the 33-vector, 14-basis KS set yields a $5 \times 9$ qutrit--qutrit BPQS (Alice has 5 inputs, Bob has 9), improving on the $7 \times 9$ strategy derived from the Peres 33-vector set (16 bases). The improvement from 7 to 5 Alice inputs arises entirely from the reduction in context count: 14 bases vs.\ 16.

This gives the island classification direct operational meaning. Each island's minimal KS set has a characteristic context count:

\begin{center}
\begin{tabular}{lcccl}
\toprule
Island & rays & bases & $\eta$ & Operational consequence \\
\midrule
Eisenstein & 33 & 14 & 0.929 & Fewest bases $\to$ simplest BPQS ($5 \times 9$) \\
Peres & 33 & 16 & 0.812 & Standard BPQS ($7 \times 9$) \\
$\Z[\sqrt{-2}]$ & 33 & 16 & 0.812 & Isomorphic to Peres \\
Integer (CK-31) & 31 & 17 & 1.000 & Fewest rays but most bases \\
Heegner-7 & 43 & 23 & 0.783 & Richest structure; unexplored BPQS \\
\bottomrule
\end{tabular}
\end{center}

\noindent The table reveals a trade-off between vector count and context count that the island framework makes explicit. CK-31 achieves the minimum number of rays (31) but at the cost of having the most bases among the 33-class islands (17 vs.\ 14). The Eisenstein island sacrifices two rays to gain a sparser basis structure, yielding a simpler Bell scenario. This trade-off---\emph{ray economy vs.\ context economy}---is invisible without the algebraic island classification.

\subsubsection{Contextuality as a quantum resource}

The CSW framework (Section~\ref{sec:csw}) quantifies contextual advantage via the graph invariants $\alpha$, $\vartheta$, and $\alpha^*$. A complementary approach, due to Abramsky, Barbosa, and Mansfield~\cite{AbramskyBarbosaMansfield2017}, defines the \emph{contextual fraction} of an empirical model as a resource monotone: it is non-increasing under classical operations and directly quantifies the quantum advantage in nonlocal games and measurement-based quantum computation. The contextual fraction is computable by linear programming and provides an operational measure of ``how contextual'' a given set of correlations is.

Our critical bases analysis (Section~\ref{sec:critical-bases}) provides a related but distinct measure at the level of the KS set itself rather than the empirical model. The essentiality ratio $\eta$ quantifies how tightly the basis hypergraph is woven: a high $\eta$ means most bases are individually necessary for the KS proof. Together with the CSW invariants and the BPQS context count, these measures provide three independent axes along which algebraic islands differ:

\begin{enumerate}
    \item \textbf{Graph-theoretic axis} (CSW): $\vartheta/\alpha$ measures the quantum--classical gap on the orthogonality graph. The Heegner-7 pool leads ($\vartheta/\alpha = 1.215$); the Peres and $\Z[\sqrt{-2}]$ pools are $\vartheta$-perfect.
    \item \textbf{Hypergraph axis} (critical bases): $\eta$ measures the rigidity of the KS proof. CK-31 leads ($\eta = 1.000$); the Heegner-7 minimized set is most slack ($\eta = 0.783$).
    \item \textbf{Operational axis} (BPQS): the context count determines Bell scenario complexity. The Eisenstein island leads (14 bases $\to$ $5 \times 9$ BPQS); CK-31 is worst among the small islands (17 bases).
\end{enumerate}

\noindent No single island dominates on all three axes. CK-31 is the tightest proof (highest $\eta$) and has a nontrivial CSW sandwich ($\alpha < \vartheta < \alpha^*$), but its 17 bases make it operationally costly. The Eisenstein island has the simplest BPQS and near-maximal tightness ($\eta = 0.929$), but only a marginal CSW advantage ($\vartheta/\alpha = 1.004$). The Heegner-7 island has the richest CSW structure but the most bases and lowest tightness. This three-way trade-off suggests that the choice of algebraic substrate has genuine consequences for quantum information tasks, and that no single notion of ``simplest'' or ``best'' KS set is appropriate---it depends on the operational criterion.

\subsubsection{Open questions}

Several questions connecting the island framework to quantum strategies remain open:

\begin{itemize}
    \item Does CK-31 serve as the engine of any perfect quantum strategy? Its 17 bases and high essentiality ratio suggest a rigid but operationally expensive BPQS.
    \item What BPQS does the Heegner-7 island's 43-vector, 23-basis set produce? Its richer orthogonality structure (107 pairs, 23 bases) may enable qualitatively different Bell scenarios---perhaps with more outputs or higher-dimensional entanglement.
    \item Can the three-way trade-off (CSW advantage, hypergraph tightness, context economy) be formalized as a Pareto frontier across algebraic islands?
    \item Is there an algebraic island that simultaneously minimizes context count \emph{and} maximizes the CSW advantage?
\end{itemize}

The fact that Cabello's ``simplest'' construction lands precisely on one of our independently discovered islands---and not between them---is strong evidence that the island classification captures genuine algebraic structure rather than an artifact of our search methodology.

\subsection{Graph-theoretic invariants and Bell inequalities}\label{sec:csw}

The Cabello--Severini--Winter (CSW) framework~\cite{CabelloSeveriniWinter2014} maps the orthogonality graph $G$ of a set of rays to a noncontextuality inequality via three invariants:
\[
\alpha(G) \leq \vartheta(G) \leq \alpha^*(G),
\]
where $\alpha(G)$ is the independence number (classical bound), $\vartheta(G)$ is the Lov\'asz theta number (quantum bound, computed by semidefinite programming), and $\alpha^*(G)$ is the fractional packing number (nonsignaling bound, computed by linear programming). A quantum advantage exists when $\vartheta(G) > \alpha(G)$.

We compute these invariants for two representations of each island: the SAT-minimized KS set (minimum cardinality) and the full ray pool (all rays generated by the alphabet before minimization). This comparison reveals a structural mechanism by which cardinality minimization destroys the CSW contextual advantage.

\begin{table}[ht]
\centering
\caption{CSW graph invariants for minimized KS sets. $\alpha$ = independence number (classical), $\vartheta$ = Lov\'asz theta (quantum, SDP via SCS with tolerance $10^{-9}$), $\alpha^*$ = fractional packing (nonsignaling, LP via HiGHS, exact). The ratio $\vartheta/\alpha$ measures the CSW contextual advantage. Equalities $\vartheta = \alpha$ are reported when $|\vartheta - \alpha| < 10^{-6}$.}
\label{tab:csw-min}
\begin{tabular}{lccccc}
\toprule
Island (minimized) & $n$ & $\alpha(G)$ & $\vartheta(G)$ & $\alpha^*(G)$ & $\vartheta/\alpha$ \\
\midrule
Integer (CK-31) & 31 & 11 & 11.71 & 12.00 & \textbf{1.065} \\
Eisenstein & 33 & 12 & 12.05 & 13.00 & 1.004 \\
Peres & 33 & 12 & 12.00 & 12.00 & 1.000 \\
$\Z[\sqrt{-2}]$ & 33 & 12 & 12.00 & 12.00 & 1.000 \\
Heegner-7 & 43 & 16 & 16.00 & 16.00 & 1.000 \\
\bottomrule
\end{tabular}
\end{table}

\begin{table}[ht]
\centering
\caption{CSW graph invariants for full (non-minimized) ray pools. $\vartheta$ computed via SCS (tolerance $10^{-9}$); $\alpha^*$ via HiGHS (exact). The ``auxiliary'' column counts rays that participate in pairwise orthogonalities but not in any complete basis (Definition~\ref{def:auxiliary}).}
\label{tab:csw-pool}
\begin{tabular}{lccccccc}
\toprule
Island (full pool) & $n$ & bases & auxiliary & $\alpha(G)$ & $\vartheta(G)$ & $\alpha^*(G)$ & $\vartheta/\alpha$ \\
\midrule
Integer & 49 & 26 & 0 (0\%) & 17 & 17.70 & 17.00 & 1.041 \\
Eisenstein & 57 & 22 & 0 (0\%) & 18 & 19.34 & 21.00 & 1.074 \\
Heegner-7 & 145 & 42 & 76 (52\%) & 46 & 55.89 & 66.00 & \textbf{1.215} \\
Peres & 49 & 16 & 16 (33\%) & 23 & 23.00 & 23.00 & 1.000 \\
$\Z[\sqrt{-2}]$ & 49 & 16 & 16 (33\%) & 23 & 23.00 & 23.00 & 1.000 \\
\bottomrule
\end{tabular}
\end{table}

\subsubsection{Auxiliary rays and the CSW advantage}

The comparison between Tables~\ref{tab:csw-min} and~\ref{tab:csw-pool} reveals a striking pattern: the Heegner-7 island's full pool has the strongest contextual advantage of any island ($\vartheta/\alpha = 1.215$), yet its minimized 43-vector set is $\vartheta$-perfect ($\vartheta = \alpha$). To explain this, we introduce a structural distinction among rays in an algebraic pool.

\begin{definition}\label{def:auxiliary}
Let $\mathcal{P}$ be a ray pool with orthogonality graph $G = (V, E)$ and set of bases (maximal orthogonal triples) $\mathcal{B}$. A ray $v \in V$ is \textbf{basis-participating} if $v$ belongs to at least one basis $B \in \mathcal{B}$, and \textbf{auxiliary} otherwise. That is, an auxiliary ray has at least one orthogonal neighbor in $G$ but does not belong to any complete orthonormal triple.
\end{definition}

Auxiliary rays contribute edges to the orthogonality graph without contributing to the hypergraph structure that determines KS-uncolorability. Since KS-uncolorability is a property of the basis hypergraph (no consistent $\{0,1\}$ coloring assigning exactly one ``1'' per basis), auxiliary rays are invisible to it. But they are fully visible to the CSW framework, which depends only on the orthogonality \emph{graph}.

\begin{observation}\label{obs:auxiliary}
\textbf{(Auxiliary rays amplify the CSW advantage.)} Among the five algebraic islands, the CSW contextual advantage $\vartheta(G)/\alpha(G)$ of the full ray pool correlates with the proportion of auxiliary rays:
\begin{center}
\begin{tabular}{lccl}
\toprule
Island & Auxiliary fraction & $\vartheta/\alpha$ (pool) & $\vartheta/\alpha$ (minimized) \\
\midrule
Heegner-7 & 52\% & \textbf{1.215} & 1.000 \\
Peres & 33\% & 1.000 & 1.000 \\
$\Z[\sqrt{-2}]$ & 33\% & 1.000 & 1.000 \\
Integer & 0\% & 1.041 & 1.065 \\
Eisenstein & 0\% & 1.074 & 1.004 \\
\bottomrule
\end{tabular}
\end{center}
Auxiliary rays constrain the independence number $\alpha(G)$ (by adding edges that prevent large independent sets) without contributing to the basis structure. The SDP relaxation $\vartheta(G)$ can exploit these additional constraints by distributing weight fractionally across the auxiliary vertices, widening the gap $\vartheta - \alpha$.
\end{observation}

The mechanism operates as follows. In the Heegner-7 pool ($n = 145$), only 69 rays participate in the 42 bases. The remaining 76 auxiliary rays add 522 orthogonal pairs to the graph, including 9 hub vertices of degree 16 (each connected to 11\% of all other vertices). These hubs severely restrict the size of independent sets: only $\alpha/n = 0.317$ of vertices can be simultaneously independent, compared to $0.469$ for the Peres and $\Z[\sqrt{-2}]$ pools (which are $\vartheta$-perfect). The SDP, however, can assign fractional weights summing to $\vartheta/n = 0.385$, yielding a 22\% quantum advantage.

When SAT minimization reduces the Heegner-7 pool to 43 vectors, it removes auxiliary rays first---they are dispensable for KS-uncolorability. But these are precisely the rays that constrained $\alpha$ and enabled the $\vartheta > \alpha$ gap. The minimized set is $\vartheta$-perfect.

\subsubsection{Further patterns}

\begin{enumerate}
    \item \textbf{Full pools consistently outperform minimized sets} for the CSW advantage. The Eisenstein pool jumps from $\vartheta/\alpha = 1.004$ (minimized) to $1.074$ (full pool). The integer pool is the exception: it drops slightly from $1.065$ to $1.041$, suggesting that CK-31 is unusually efficient at preserving CSW structure despite minimization. Notably, the integer and Eisenstein pools have zero auxiliary rays---every ray participates in a basis---so their CSW advantages arise from graph regularity rather than auxiliary edges.

    \item \textbf{The nonsignaling bound is tightly constrained by basis structure.} For the four minimized sets where $\vartheta = \alpha$ (Peres, $\Z[\sqrt{-2}]$, Heegner-7, plus Eisenstein near-equality), we also have $\alpha^* = \alpha$: all three invariants collapse. Only the integer island (CK-31) maintains a nontrivial sandwich $\alpha = 11 < \vartheta = 11.71 < \alpha^* = 12$. In the full pools, the Heegner-7 pool shows the widest separation: $\alpha = 46 < \vartheta = 55.89 < \alpha^* = 66$, with the nonsignaling bound 43\% above the classical bound.

    \item \textbf{KS-uncolorability does not imply CSW violation}: The Peres and $\Z[\sqrt{-2}]$ islands remain $\vartheta$-perfect even in their full pools, despite having 16 auxiliary rays each. The auxiliary ray mechanism is necessary but not sufficient; the specific placement of orthogonalities matters. In these two pools, the auxiliary rays have low degree (minimum degree 3), insufficient to constrain $\alpha$ below the SDP threshold.
\end{enumerate}

\begin{remark}
The distinction between basis-participating and auxiliary rays reveals that KS-uncolorability and CSW contextual advantage are driven by different structural features of the same algebraic pool. Uncolorability requires bases arranged so that no valid $\{0,1\}$ coloring exists---a property of the basis \emph{hypergraph}. The CSW advantage requires pairwise orthogonalities arranged so that $\vartheta > \alpha$---a property of the orthogonality \emph{graph}. Auxiliary rays contribute to the latter without affecting the former. The operational lesson is that the contextual strength of an algebraic island, in the CSW sense, is determined by its full ray pool, not its minimal KS subset. Minimizing for cardinality can destroy the CSW advantage entirely.
\end{remark}

\subsubsection{Spectral explanation of the CSW advantage}\label{sec:spectral}

The auxiliary ray mechanism explains \emph{which} structural feature drives the CSW advantage. We now ask \emph{why} certain pools have $\vartheta > \alpha$ while others do not, using spectral graph theory.

The Hoffman bound provides a spectral explanation for which pools admit a CSW violation. For any graph $G$ with adjacency matrix $A$ having largest eigenvalue $\lambda_{\max}$ and smallest eigenvalue $\lambda_{\min}$, the independence number satisfies~\cite{Haemers1995}
\[
\alpha(G) \leq \frac{n \cdot (-\lambda_{\min})}{\lambda_{\max} - \lambda_{\min}}.
\]
This bound is valid for all graphs (not only regular graphs); for regular graphs $\lambda_{\max}$ equals the common degree, but the inequality holds in general with $\lambda_{\max}$ as the spectral radius. However, for disconnected graphs the bound can be loose, since the spectral radius of each component contributes independently. When $\alpha(G)$ is well below this bound, the graph is ``spectrally tight'': the eigenvalue structure constrains independent sets more than the combinatorial structure alone requires, leaving room for the SDP relaxation $\vartheta(G)$ to exceed $\alpha(G)$.

\begin{table}[ht]
\centering
\caption{Spectral properties of full ray pool orthogonality graphs. $H$ = Hoffman bound on $\alpha$. The ratio $\alpha/H$ measures spectral tightness: values near or below 1 correlate with CSW violations.}
\label{tab:spectral}
\begin{tabular}{lccccccc}
\toprule
Pool & $n$ & $\lambda_{\max}$ & $\lambda_{\min}$ & $H$ & $\alpha$ & $\alpha/H$ & $\vartheta/\alpha$ \\
\midrule
Heegner-7 & 145 & 8.25 & $-4.50$ & 51.2 & 46 & \textbf{0.898} & \textbf{1.215} \\
Eisenstein & 57 & 6.42 & $-3.45$ & 19.9 & 18 & \textbf{0.904} & \textbf{1.074} \\
Integer & 49 & 5.98 & $-3.29$ & 17.4 & 17 & 0.977 & 1.041 \\
Peres & 49 & 5.53 & $-3.68$ & 19.6 & 23 & 1.176 & 1.000 \\
$\Z[\sqrt{-2}]$ & 49 & 5.53 & $-3.68$ & 19.6 & 23 & 1.176 & 1.000 \\
\bottomrule
\end{tabular}
\end{table}

The pattern in Table~\ref{tab:spectral} is monotonic: the three pools with $\alpha/H < 1$ (spectrally tight) have CSW violations, with the largest violation ($\vartheta/\alpha = 1.215$) at the tightest ratio ($\alpha/H = 0.898$). The two $\vartheta$-perfect pools have $\alpha/H > 1$, meaning $\alpha$ exceeds the Hoffman bound; this occurs because these graphs are disconnected (the 16 auxiliary rays of degree 3 form loosely connected components), invalidating the bound for non-connected graphs. The spectral gap $\lambda_{\max} - \lambda_2$ also correlates: the Heegner-7 pool has the largest spectral gap (4.09), consistent with graph-theoretic results linking expansion to large $\vartheta/\alpha$ ratios.

\subsection{Contextual tightness: critical bases analysis}\label{sec:critical-bases}

The CSW framework measures how much quantum mechanics exceeds classical bounds on the orthogonality \emph{graph}. We now introduce a complementary measure that probes the internal structure of the basis \emph{hypergraph}: how many bases must be removed before a KS set becomes colorable?

\begin{definition}\label{def:essential-basis}
Let $\mathcal{S}$ be a KS set with basis set $\mathcal{B} = \{B_1, \ldots, B_m\}$. A basis $B_i$ is \textbf{essential} if $\mathcal{S}$ with the constraint from $B_i$ removed admits a valid $\{0,1\}$ coloring---i.e., the set is KS-uncolorable only when $B_i$'s ``exactly one green'' rule is enforced. The \textbf{critical number} $\kappa(\mathcal{S})$ is the minimum $k$ such that some $k$-element subset of bases can be removed to make $\mathcal{S}$ colorable. The \textbf{contextual fraction} is $\kappa(\mathcal{S}) / |\mathcal{B}|$.
\end{definition}

Note that removing a basis relaxes the ``exactly one green per triple'' constraint for that context, while all pairwise orthogonality constraints (``at most one green'') remain---orthogonality is a geometric fact, not a contextual choice. A set with $\kappa = 1$ is \emph{fragile}: a single basis removal suffices to restore colorability. A set with large $\kappa$ is \emph{robust}: many independent constraints must be simultaneously relaxed.

We compute these quantities for each island's SAT-minimized KS set and, for comparison, the full ray pools of the Peres and Heegner-7 islands. The \textbf{essentiality ratio} $\eta = (\text{essential bases}) / |\mathcal{B}|$ measures what fraction of bases are individually indispensable.

\begin{table}[ht]
\centering
\caption{Critical bases analysis. $|\mathcal{B}|$ = number of bases, ess = essential bases (whose individual removal breaks uncolorability), $\eta$ = essentiality ratio, $\kappa$ = critical number (minimum bases to remove), CF = contextual fraction $\kappa/|\mathcal{B}|$.}
\label{tab:critical-bases}
\begin{tabular}{lcccccc}
\toprule
Set & rays & $|\mathcal{B}|$ & ess & $\eta$ & $\kappa$ & CF \\
\midrule
Integer (CK-31) & 31 & 17 & \textbf{17} & \textbf{1.000} & 1 & 0.059 \\
Eisenstein min-33 & 33 & 14 & 13 & 0.929 & 1 & 0.071 \\
Peres min-33 & 33 & 16 & 13 & 0.812 & 1 & 0.063 \\
$\Z[\sqrt{-2}]$ min-33 & 33 & 16 & 13 & 0.812 & 1 & 0.063 \\
Heegner-7 min-43 & 43 & 23 & 18 & 0.783 & 1 & 0.044 \\
\midrule
Peres pool & 49 & 16 & 13 & 0.812 & 1 & 0.063 \\
Heegner-7 pool & 145 & 42 & 0 & 0.000 & \textbf{2} & 0.048 \\
\bottomrule
\end{tabular}
\end{table}

\subsubsection{Structure of the results}

Three findings emerge from Table~\ref{tab:critical-bases}.

\begin{enumerate}
    \item \textbf{CK-31 is maximally tight.} Every one of its 17 bases is essential: $\eta = 1.000$. No basis is redundant for the KS proof. This is a structural consequence of CK-31's minimality---at 31 vectors, it is the smallest known KS set in $\R^3$, and it achieves this by having every constraint contribute non-redundantly to the coloring obstruction. The integer island's KS proof is a perfectly rigid structure with no slack.

    \item \textbf{Peres and $\Z[\sqrt{-2}]$ are again identical.} Both have 16 bases, 13 essential, 3 removable, and the same essentiality ratio of 0.812. This extends the graph isomorphism of Proposition~\ref{prop:isomorphism} to the hypergraph level: not only do these two islands have isomorphic orthogonality graphs, they have identical critical bases structure. The three removable bases occupy the same structural positions in both sets.

    \item \textbf{The Heegner-7 full pool is the only set with $\kappa > 1$.} No single basis removal breaks uncolorability---every basis is individually dispensable (all 42 bases are removable, $\eta = 0$). But removing certain \emph{pairs} of bases does break it: we find exactly 24 critical pairs out of $\binom{42}{2} = 861$ possible pairs, a critical pair density of $0.028$. This high redundancy---every basis is backed up by others---reflects the Heegner-7 pool's large size (145 rays, 42 bases) compared to the other islands.
\end{enumerate}

\subsubsection{Relationship to other invariants}

The essentiality ratio $\eta$ provides a measure of contextual tightness that is independent of both the CSW advantage $\vartheta/\alpha$ and the set cardinality:

\begin{center}
\begin{tabular}{lcccc}
\toprule
Island (minimized) & rays & $\eta$ & $\vartheta/\alpha$ & bases/ray \\
\midrule
Integer (CK-31) & 31 & 1.000 & 1.065 & 0.55 \\
Eisenstein & 33 & 0.929 & 1.004 & 0.42 \\
Peres & 33 & 0.812 & 1.000 & 0.48 \\
$\Z[\sqrt{-2}]$ & 33 & 0.812 & 1.000 & 0.48 \\
Heegner-7 & 43 & 0.783 & 1.000 & 0.53 \\
\bottomrule
\end{tabular}
\end{center}

CK-31 ranks first in both $\eta$ and CSW advantage among the minimized sets, but the two measures are not monotonically related: the Eisenstein set has $\eta = 0.929$ (second-highest) but $\vartheta/\alpha = 1.004$ (second-lowest among those with any advantage at all). This confirms that $\eta$ captures a genuinely different aspect of contextual structure---the tightness of the hypergraph proof rather than the graph-theoretic gap between classical and quantum bounds.

The critical number $\kappa = 1$ for all minimized sets means that SAT minimization produces sets where at least one basis is a ``bottleneck'' for the proof. This is expected: minimization removes redundant rays (and hence bases) until no further removal is possible, leaving a structure where some bases are load-bearing. The transition to $\kappa = 2$ in the Heegner-7 full pool shows that algebraic completeness provides a qualitatively different kind of robustness---contextual redundancy---that minimization destroys.

\subsection{Limitations and open questions}

Our search is computational, not exhaustive over all possible number fields. Specifically:

\begin{enumerate}
    \item \textbf{Alphabet choice within fields}: For $d \equiv 1 \pmod{4}$, the ring of integers of $\Q(\sqrt{d})$ is $\Z[(1+\sqrt{d})/2]$, not $\Z[\sqrt{d}]$. Our discovery that $\varphi = (1+\sqrt{5})/2$ produces a KS set while $\sqrt{5}$ does not shows that the choice of generator matters. We have tested the ring-of-integers generator $(1+\sqrt{d})/2$ for $d \in \{5, 13, 17, 21, 29\}$; results are reported in the supplementary data.
    \item \textbf{Complex quadratic fields}: we have exhaustively tested all nine imaginary quadratic fields with class number 1 (the Heegner numbers), as well as $\Z[\sqrt{-d}]$ for $d = 2, 5, 6, 7, 10, 11, 13, 14, 15$ with enriched alphabets and Hermitian completion. The classification appears complete for imaginary quadratic fields.
    \item \textbf{Cubic and higher extensions}: we have tested cubic extensions $\Q(\sqrt[3]{d})$ for $d = 2, 3, 4, 5$ (all colorable). Higher-degree algebraic extensions remain unexplored.
    \item \textbf{Completion termination}: the cross-product completion operation (Section~\ref{sec:closure}) terminates empirically for all tested alphabets, but we do not prove termination in general. For alphabets whose coordinate ring is not closed under the cofactor formula's operations, the completed pool could in principle grow without bound. In practice, we impose a maximum iteration limit and verify stabilization.
    \item \textbf{Transcendental coordinates}: our methods require algebraic coordinates (for exact orthogonality testing); KS sets with transcendental coordinates are not excluded.
    \item \textbf{The $6 \mid n$ conjecture}: verified for $n \leq 30$ but not proved for all $n$. A proof would likely follow from the representation theory of cyclic groups acting on $\C^3$.
    \item \textbf{Minimization is heuristic, but 33 and 43 appear optimal}: our randomized greedy minimization does not certify optimality in general. However, for the three 33-vector islands (Peres, Eisenstein, $\Z[\sqrt{-2}]$), we have verified that \emph{no} 32-vector KS subset exists within any known minimal 33-set: across 2{,}000 trials per island, we found 1--3 distinct 33-sets each, and tested all possible single-vector removals (33--99 subsets of size 32 per island)---all were colorable. Similarly, for the Heegner-7 island, all 258 subsets of size 42 across 6 distinct 43-sets were colorable. While this does not exclude the possibility of a 32-vector set that is not a subset of any known 33-set, it strongly suggests that 33 and 43 are optimal for these pools. The integer island's minimum of 31 is supported by exhaustive enumeration over the 37-ray union of discovered 31-sets (Proposition~\ref{prop:31optimal}); a certificate-based approach covering the full pool is left to future work.
    \item \textbf{Lower bound improvement}: our work does not improve the lower bound of 24; it constrains where a sub-31 construction could come from.
    \item \textbf{Operational significance}: our connection to Cabello's work (Section~\ref{sec:cabello}) is primarily observational---we identify the Eisenstein island with his simplest KS set, but do not prove that other islands cannot arise as engines of perfect quantum strategies. A systematic study of which Bell scenarios' perfect strategies correspond to which algebraic islands would deepen this connection.
\end{enumerate}

\section{Conclusion}\label{sec:conclusion}

We have presented a systematic computational investigation of the algebraic landscape underlying Kochen--Specker sets in three dimensions. The principal findings are:

\begin{enumerate}
    \item Among the number fields and alphabets surveyed, KS sets in $\R^3$ and $\C^3$ cluster into exactly six algebraic islands, with smallest KS subsets found at $31 < 33 = 33 = 33 < 43 < 52$ vectors, corresponding to five distinct cancellation mechanisms. The Peres, complex quadratic, and Gaussian islands share the same ``norm-2'' identity and produce \emph{graph-isomorphic} orthogonality graphs at both full pool and minimized levels (Proposition~\ref{prop:isomorphism}), reducing six algebraic islands to five distinct graph types.
    \item The golden ratio field $\Q(\varphi)$, the complex quadratic ring $\Z[\sqrt{-2}]$, and the Heegner-7 ring $\Z[(1+\sqrt{-7})/2]$ are newly discovered islands. The golden ratio island is invisible to raw alphabet searches and revealed only by cross-product completion; the complex quadratic island is the complex analogue of the Peres island with isomorphic combinatorial structure; the Heegner-7 island fills the gap between the 33-cluster and the Golden island at a novel minimum of 43 vectors.
    \item A systematic survey of all nine imaginary quadratic fields with class number 1 (the Heegner numbers $d = 1, 2, 3, 7, 11, 19, 43, 67, 163$) shows that KS-uncolorability requires the ring-of-integers generator to have Hermitian norm $\leq 2$. Only $d = 1, 2, 3, 7$ satisfy this bound.
    \item Strong computational evidence supports 31 as the minimum KS set size for the integer island: all $10.3 \times 10^6$ subsets of the 37 rays appearing in discovered minimal 31-sets are colorable (Proposition~\ref{prop:31optimal}).
    \item The $n$th roots of unity produce KS-uncolorable sets if and only if $6 \mid n$ (verified for $n \leq 30$), connecting cyclotomic number theory to the $\Z[1/6]$ minimality result of Cortez, Morales, and Reyes.
    \item The realizability gap is vast and uniform: across 30{,}000+ abstract hypergraph trials at sizes 22--30, over 10{,}000 were KS-uncolorable but none were realizable in $\R^3$. The geometry of three-dimensional space is the bottleneck, not combinatorics.
    \item The Eisenstein island's 33-vector KS set is identified with Cabello's ``simplest KS set,'' connecting the algebraic island framework to the proof that KS sets are the engine of bipartite perfect quantum strategies.
    \item The Conway--Kochen 31-vector set is not robust under coordinate perturbation: 1\% noise destroys all exact orthogonalities, consistent with the rigidity result of Trandafir and Cabello.
    \item The Cabello--Severini--Winter graph invariants reveal that \emph{auxiliary rays}---rays that participate in pairwise orthogonalities but not in any complete basis (Definition~\ref{def:auxiliary})---amplify the quantum contextual advantage (Observation~\ref{obs:auxiliary}). The Heegner-7 full pool, with 52\% auxiliary rays, achieves $\vartheta/\alpha = 1.215$ (a 22\% quantum advantage), the strongest of any island. Its minimized 43-vector set, stripped of auxiliary rays, has $\vartheta = \alpha$. This demonstrates that KS-uncolorability (a property of the basis hypergraph) and CSW violation (a property of the orthogonality graph) are structurally independent: auxiliary rays contribute to the latter without affecting the former.
    \item The KS theorem is a rounding impossibility: the failure of consistently rounding Born-rule probabilities to deterministic $\{0, 1\}$ valuations. The 24-to-31 gap measures the cost geometric realizability imposes on abstract combinatorial constraints.
    \item A critical bases analysis (Section~\ref{sec:critical-bases}) reveals that CK-31 is maximally tight: all 17 of its bases are individually essential for the KS proof ($\eta = 1.000$), while the other minimized islands have essentiality ratios ranging from 0.783 (Heegner-7) to 0.929 (Eisenstein). All minimized sets have critical number $\kappa = 1$, but the Heegner-7 full pool achieves $\kappa = 2$---no single basis removal breaks its uncolorability---demonstrating that algebraic completeness provides contextual redundancy that minimization destroys. The essentiality ratio $\eta$ is independent of both the CSW advantage and set cardinality, capturing a distinct aspect of contextual structure: the tightness of the hypergraph proof.
\end{enumerate}

\subsection*{Reproducibility}

All code and data are available at \url{https://github.com/michaelkernaghan/contextuality}. The repository includes explicit ray coordinates and triad/basis lists for all islands, including the 43-vector Heegner-7 set and the 52-vector golden-ratio set, as well as verifier scripts that reconstruct each KS set from its alphabet and confirm UNSAT via SAT solving. The implementation uses Python 3.11 with PySAT 1.8 (Glucose4 solver), NumPy, SciPy (L-BFGS-B optimizer and HiGHS LP solver), and CVXPY (SCS semidefinite programming solver, tolerance $10^{-9}$). Tables~\ref{tab:alphabets}--\ref{tab:closure} are generated by \texttt{ks\_islands.py}; Table~\ref{tab:roots} by \texttt{ks\_complex.py}; Tables~\ref{tab:realizability}--\ref{tab:realizability-large} by \texttt{ks\_geometry.py}; Table~\ref{tab:complex-quadratic} and the exhaustive 31-optimality check by \texttt{ks\_new\_islands.py} and \texttt{ks\_sat.py}; Table~\ref{tab:heegner} by \texttt{ks\_new\_island\_analysis.py}; Tables~\ref{tab:csw-min}--\ref{tab:csw-pool} and~\ref{tab:spectral} by \texttt{ks\_csw.py}, \texttt{ks\_csw\_extended.py}, and \texttt{ks\_graph\_analysis.py}; Table~\ref{tab:critical-bases} by \texttt{ks\_critical\_bases.py}; Proposition~\ref{prop:isomorphism} by \texttt{ks\_graph\_analysis.py}. All randomized experiments use \texttt{random.seed(42)} for reproducibility.

\subsection*{Acknowledgments}

The author thanks Ad\'an Cabello for pointing out the connection between the Eisenstein island and his ``simplest KS set'' construction. The author acknowledges the use of Claude (Anthropic) for computational assistance, literature synthesis, and manuscript preparation.

\begin{thebibliography}{99}

\bibitem{KochenSpecker1967}
S.~Kochen and E.~P.~Specker,
``The problem of hidden variables in quantum mechanics,''
\textit{J. Math. Mech.} \textbf{17}, 59--87 (1967).

\bibitem{Peres1991}
A.~Peres,
``Two simple proofs of the Kochen--Specker theorem,''
\textit{J. Phys. A: Math. Gen.} \textbf{24}, L175--L178 (1991).

\bibitem{Peres1993}
A.~Peres,
\textit{Quantum Theory: Concepts and Methods}
(Kluwer Academic Publishers, 1993).

\bibitem{Kernaghan1994}
M.~Kernaghan,
``Bell--Kochen--Specker theorem for 20 vectors,''
\textit{J. Phys. A: Math. Gen.} \textbf{27}, L829--L830 (1994).

\bibitem{KernaghanPeres1995}
M.~Kernaghan and A.~Peres,
``Kochen--Specker theorem for eight-dimensional space,''
\textit{Phys. Lett. A} \textbf{198}, 1--5 (1995).

\bibitem{CabelloEstebaranzGarcia1996}
A.~Cabello, J.~M.~Estebaranz, and G.~Garc\'{\i}a-Alcaine,
``Bell--Kochen--Specker theorem: A proof with 18 vectors,''
\textit{Phys. Lett. A} \textbf{212}, 183--187 (1996).

\bibitem{Pavicic2004}
M.~Pavicic, J.-P.~Merlet, B.~McKay, and N.~D.~Megill,
``Kochen--Specker vectors,''
\textit{J. Phys. A: Math. Gen.} \textbf{38}, 1577--1592 (2005).
arXiv:quant-ph/0409014.

\bibitem{CabelloSeveriniWinter2014}
A.~Cabello, S.~Severini, and A.~Winter,
``Graph-theoretic approach to quantum correlations,''
\textit{Phys. Rev. Lett.} \textbf{112}, 040401 (2014).

\bibitem{UijlenWesterbaan2016}
S.~Uijlen and B.~Westerbaan,
``A Kochen--Specker system has at least 22 vectors,''
\textit{New Generation Computing} \textbf{34}, 3--23 (2016).

\bibitem{AudemardSimon2018}
G.~Audemard and L.~Simon,
``On the Glucose SAT solver,''
\textit{Int. J. Artif. Intell. Tools} \textbf{27}, 1840001 (2018).

\bibitem{KunjwalSpekkens2018}
R.~Kunjwal and R.~W.~Spekkens,
``Beyond the Cabello--Severini--Winter framework: Making sense of contextuality without sharpness of measurements,''
\textit{Phys. Rev. A} \textbf{97}, 052110 (2018).

\bibitem{CortezMoralesReyes2022}
V.~Cortez, R.~Morales, and E.~J.~Reyes,
``Minimal ring extensions of the integers exhibiting Kochen--Specker contextuality,''
arXiv:2211.13216 (2022).

\bibitem{LiBrightGanesh2024}
J.~Li, C.~Bright, and V.~Ganesh,
``Leveraging SAT solvers for the minimum Kochen--Specker problem,''
in \textit{Proceedings of the 33rd International Joint Conference on Artificial Intelligence (IJCAI-24)}, pp.~1944--1952 (2024).
arXiv:2306.13319.

\bibitem{Cabello2024bipartite}
A.~Cabello,
``Every bipartite perfect quantum strategy requires a Kochen--Specker set,''
\textit{Phys.\ Rev.\ Lett.} \textbf{134}, 010201 (2025).
arXiv:2311.17735.

\bibitem{Cabello2024equivalences}
A.~Cabello,
``Bipartite perfect quantum strategies, full Bell nonlocality, algebraic proofs of the KS theorem, and faces of the nonsignaling polytope are equivalent,''
\textit{Phys.\ Rev.\ Research} \textbf{6}, L042035 (2024).
arXiv:2310.10600.

\bibitem{Cabello2025simplest}
A.~Cabello,
``Simplest Kochen--Specker set,''
\textit{Phys.\ Rev.\ Lett.} (2025).
arXiv:2508.07335.

\bibitem{TrandafirCabello2025}
S.~Trandafir and A.~Cabello,
``Two fundamental solutions to the rigid Kochen--Specker set problem and the solution to the minimal Kochen--Specker set problem under one assumption,''
\textit{Phys.\ Rev.\ A} \textbf{111}, 052204 (2025).
arXiv:2501.11640.

\bibitem{TrandafirCabello2025optimal}
S.~Trandafir and A.~Cabello,
``Optimal conversion of Kochen--Specker sets into bipartite perfect quantum strategies,''
\textit{Phys.\ Rev.\ A} \textbf{111}, 022408 (2025).
arXiv:2410.17470.

\bibitem{AbramskyBarbosaMansfield2017}
S.~Abramsky, R.~S.~Barbosa, and S.~Mansfield,
``The contextual fraction as a measure of contextuality,''
\textit{Phys.\ Rev.\ Lett.} \textbf{119}, 050504 (2017).
arXiv:1705.07918.

\bibitem{Vazirani2001}
V.~V.~Vazirani,
\textit{Approximation Algorithms}
(Springer, 2001).

\bibitem{Gleason1957}
A.~M.~Gleason,
``Measures on the closed subspaces of a Hilbert space,''
\textit{J. Math. Mech.} \textbf{6}, 885--893 (1957).

\bibitem{Haemers1995}
W.~H.~Haemers,
``Interlacing eigenvalues and graphs,''
\textit{Linear Algebra Appl.} \textbf{226--228}, 593--616 (1995).

\end{thebibliography}

\end{document}
